% GCT_POS branch /root/nonstandard. Standalone section, intended for \input.
% Model: GPT-6 Astra Pro; effort setting not exposed. 2026-11-21 UTC.
\section{Nonstandard positivity, integral orders, and dimension rings}
\label{sec:nonstandard-new}

\paragraph{Scope and status.}
N-009 settles the integral basis question negatively: the actual rank-four order is not locally free at $(2,v-1)$ and has no $\mathbb Z[v,v^{-1}]$-basis. The section also gives N-001--N-008 and N-010 with their stated scope. N-010 supplies a complete rational-$D$ split model and GT path data, without identifying them with the conjectural canonical globalization. N-012 disproves the literal variable-scaling minimum in VIII's external recipe. A bar-fixed repair is a distinct problem; the section does not claim a designated $E(4)$ table from that repair.  N-002 and N-003 have exact finite
certificates, independently replayed through standard Hecke multiplication.
N-004 contains an all-degree conditional theorem and an unconditional
rational-coefficient corollary.  N-005 has a direct algebraic proof as well
as exhaustive exact checks of its finite multiplication tables.

Put $A=\mathbb Z[v,v^{-1}]$, $K=\mathbb Q(v)$,
$\delta=v+v^{-1}$ and $D=\delta^2$.  The standard Hecke convention is
$(T_s-v)(T_s+v^{-1})=0$, with $C'_s=T_s+v^{-1}$ and $C_s=T_s-v$.
Let $\mathcal B_r$ be the unital $A$-subalgebra of
$H_r\otimes_A H_r$ generated by
\[
 P_i=C'_{s_i}\otimes C'_{s_i}+C_{s_i}\otimes C_{s_i}.
\]
The cone in question is
$\bigoplus_{a,b\in S_r}\mathbb N[v,v^{-1}](C_a\otimes C_b)$.
The generic dimension $\dim_K K\mathcal B_4=114$ is an inherited
literature result, not a discovery here: GCT IV, subsection 11.8, printed p.~67,
``The algebra $\check H_4$,'' following Remark 11.19; its nine simple
module dimensions are $1,2,3,6,5,5,2,3,1$.
All permutation coordinates below are one-line notation.

The integral object is precisely GCT IV, Definition 11.1, equation (93),
printed p.~59, with the ground ring fixed in its subsection 2.1, p.~16.
No localization or saturation enters that definition.  The neighboring
localization $A[1/[2]]$ is used for normalized projectors.  The canonical
generator convention above is fixed in GCT IV, equation (22), p.~20,
and subsection 3.1, p.~21, and independently in Blasiak (2015),
Definition 2.2 and its preceding conventions.  It resolves an isolated
reversal of the parameter powers in GCT IV's introductory paragraph on
p.~58.  See the accompanying definition/literature audit for the exact
source comparison and the bounded historical search.

\subsection{N-009: the full integral rank-four order is not locally free}
\label{sec:nonstandard-nonfree}

\paragraph{Theorem.}
For the integral order $\mathcal B_4\subset H_4\otimes_A H_4$ generated
by the actual $P_1,P_2,P_3$, with $A=\mathbb Z[v,v^{-1}]$, one has
\[
 \dim_{\mathbb F_2}\mathcal B_4/(2,v-1)\mathcal B_4\ \ge126>114
       =\operatorname{rank}_A\mathcal B_4.
 \tag{N-16}
\]
Consequently $\mathcal B_4$ is not locally free at $(2,v-1)$ and has
no $A$-basis.  In particular it has no $A$-basis in the $C\otimes C$
cone or in any signed or factored variant of that cone, in any degree.
The conclusion concerns this specified integral order; it does not
exclude a basis over $\mathbb Q[v,v^{-1}]$, or a different free
integral lattice in the same generic algebra.

\paragraph{The finite ring and the all-degree implication.}
Let
\[
 R=(\mathbb Z/2^{10}\mathbb Z)[\eta]/(\eta^8),\qquad
 v\longmapsto1+\eta,
 \qquad\mathfrak m=(2,\eta)\subset R.
 \tag{N-17}
\]
The image of $v$ is a unit, with inverse
$\sum_{j=0}^7(-\eta)^j$, so (N-17) defines a homomorphism $A\to R$.
Let $M$ be the image of $\mathcal B_4$ in
$(H_4\otimes_A H_4)\otimes_A R$.  There is a surjection of
$R$-modules
\[
                  R\otimes_A\mathcal B_4\twoheadrightarrow M.
 \tag{N-18}
\]
Neither injectivity nor flat base change is assumed.  Passing to the
residue field gives a surjection
\[
 \mathcal B_4/(2,v-1)\mathcal B_4\twoheadrightarrow M/\mathfrak mM.
 \tag{N-19}
\]
The exact finite calculation below proves
\[
 |M|=2^{3690},\qquad|\mathfrak mM|=2^{3564},\qquad
 \dim_{\mathbb F_2}M/\mathfrak mM=126.
 \tag{N-20}
\]
This proves (N-16).  The order is finite over $A$, since it is a
submodule of a finite free module over the Noetherian ring $A$.
It is torsion-free and has generic rank 114, the dimension established
in GCT IV, subsection 11.8, printed p.~67, ``The algebra $\check H_4$.''
If its localization at $(2,v-1)$ were free, its rank would therefore
be 114 and its residue space would have dimension 114, contradicting
(N-16).  Equivalently, an $A$-basis would yield a surjection
$R^{114}\twoheadrightarrow M$, impossible by (N-20) or Nakayama's
lemma.  Thus a finite closed image proves the assertion for the full
order in every degree; it is not an inference from a word-length cutoff.

\paragraph{Coordinate conventions.}
List the 24 permutations lexicographically, in zero-based one-line
notation.  Use the 300 basis vectors of the flip-invariant ambient
space consisting of $T_a\otimes T_a$ and
$T_a\otimes T_b+T_b\otimes T_a$ for $a<b$.
An off-diagonal coordinate is the coefficient of either ordered
$T$-tensor, with no division by two.  This remains a free basis in
characteristic two: the two ordered tensor coordinates are distinct.
If $j\in\{0,\ldots,299\}$ is the coordinate index, the scalar
coefficient of $\eta^h$ is stored at index $j+300h$.
There are 2400 additive coordinates in $\mathbb Z/1024\mathbb Z$.
All $P$-words are flip-invariant, so this coordinate space contains
$M$ faithfully.

Right multiplication is computed using
\[
 \begin{split}
 T_wT_{s_i}&=T_{ws_i}+
     \mathbf1_{\ell(ws_i)<\ell(w)}(v-v^{-1})T_w,\\
 P_i&=2T_{s_i}\otimes T_{s_i}
 -(v-v^{-1})(T_{s_i}\otimes1+1\otimes T_{s_i})
 +(v^2+v^{-2})(1\otimes1).
 \end{split}
 \tag{N-21}
\]
These formulas use integer arithmetic in (N-17).

\paragraph{The finite generating certificate and closure.}
The file \nolinkurl{mixed_jet_p2_k10_L8.json} lists 128 actual
$P$-words.  It includes the empty word and has the following counts:
\[
\begin{array}{c|rrrrrrrrrr}
\text{length}&0&1&2&3&4&5&6&7&8&9\\\hline
\text{accepted words}&1&3&5&8&13&19&26&33&19&1.
\end{array}
\]
Let $g_w$ denote each specialized word vector.  The certificate spans
$M_0=\sum_w Rg_w$; as an additive module it is generated by all
$\eta^h g_w$, $0\le h<8$.
The exact verification checks all $128\cdot3=384$ inclusions
$g_wP_i\in M_0$.  Since scalar multiplication commutes with the
$P_i$, this proves $M_0P_i\subseteq M_0$.  Also $1\in M_0$ and
$M_0$ is an $R$-module.  Induction on arbitrary word length gives
$M\subseteq M_0$, while each listed word belongs to the actual image,
so $M_0\subseteq M$.  Hence $M=M_0$.

The additive module $N=2M+\eta M$ is generated exactly by
\[
 \{2g_w:w\}\ \cup\
 \{\eta^h g_w:w,\ 1\le h<8\}.
 \tag{N-22}
\]
In particular this is multiplication by the maximal ideal inside
$M$, not intersection with an ambient maximal-ideal submodule.
Confusing those two operations would destroy the obstruction.

\paragraph{Cardinality lemma for the triangular certificate.}
For $R_0=\mathbb Z/p^k\mathbb Z$, suppose rows $r_1,\ldots,r_t$
have successively increasing first nonzero columns $j_i$, normalized
entries $(r_i)_{j_i}=p^{a_i}$, $0\le a_i<k$, and
\[
                 p^{k-a_i}r_i\in
                   \operatorname{span}_{R_0}(r_{i+1},\ldots,r_t).
 \tag{N-23}
\]
Then the cardinality of their additive span is
$p^{\sum_i(k-a_i)}$.  To prove this, project the span of rows
$i,\ldots,t$ onto column $j_i$.  Its image is the ideal
$p^{a_i}R_0$, of size $p^{k-a_i}$.  A linear combination in the kernel
has coefficient of $r_i$ divisible by $p^{k-a_i}$; (N-23) then puts
it in the span of the later rows.  The resulting short exact sequence
and induction prove the formula.  Standard leading-column reduction
also decides membership: if the current leading coordinate is not
divisible by its normalized pivot, membership fails; otherwise
subtract that multiple and continue.  Different choices of the
multiple differ by (N-23), so this reduction is valid over $R_0$.

The search constructs such a basis by exact row insertion.  A pivot
is normalized using an invertible scalar; when a new pivot has
smaller valuation it replaces the old pivot, which is reinserted.
The row $p^{k-a_i}r_i$ is inserted as well.  Each genuine enlargement
strictly increases the finite additive subgroup, so the process
terminates.  The stored rows in
\nolinkurl{mixed_jet_p2_k10_L8_rows.json.gz} satisfy all the
normalization and annihilator conditions (N-23), and all generator
membership and closure tests.  Their cardinality exponents are 3690
and 3564.

\paragraph{Independent local Smith verification.}
A second algorithm works directly from the 128 word vectors, without
using triangular row insertion.  Among the remaining entries choose
one of globally smallest $p$-valuation $a$.  Normalize it to $p^a$
by a unit row scaling, and eliminate that column in all later rows.
Every coordinate of the pivot row is divisible by $p^a$, so its
additive order is exactly $p^{k-a}$ and its annihilator is the zero
row.  Previously cleared pivot columns stay zero.  The same projection
argument therefore multiplies the span cardinality by $p^{k-a}$ at
each pivot; a zero residual block finishes the calculation.

For $M$ and $N=\mathfrak mM$, the numbers of cyclic factors with
valuation $a=0,\ldots,9$ in this independent computation are
\[
\begin{array}{c|rrrrrrrrrr}
a&0&1&2&3&4&5&6&7&8&9\\\hline
M&62&46&78&90&106&62&32&46&64&62\\
N&45&54&72&90&99&71&34&46&70&59.
\end{array}
\tag{N-24}
\]
Thus $\sum_a(10-a)n_a$ equals 3690 and 3564 respectively, proving
(N-20).  The quotient is killed by 2 and by $\eta$, so its order
$2^{126}$ is exactly its $\mathbb F_2$-dimension exponent.

The parent agent's independent program
\nolinkurl{branches/root/verify_mixed_jet.py} imports none of this
branch's code.  It implements (N-21) directly, checks all words and
closure, and uses the global-valuation Smith method just described.
Its output is
\nolinkurl{reviews/mixed_jet_p2_independent_verification.json}.
A third replay, \nolinkurl{verify_mixed_jet_certificate.py}, first
constructs words as full Laurent polynomials and only then specializes
to (N-17), using iterated unit elimination and division by 2 for the
cardinality calculation.  It yields exactly the same (N-24), and
also checks the stored triangular certificate.  Small-ring exhaustive
membership/cardinality tests of the row-insertion routine passed.
As a calibration, the identical mixed jet test at rank three gives
$\dim M/\mathfrak mM=10$, consistent with the known rank-three free
order.  The odd-prime checks $(p,k,L)=(3,2,8),(3,3,12)$ give 114 and
are inconclusive; they are not used to prove freeness anywhere.

\paragraph{Relation to N-003 and N-004.}
The explicit positive generic basis and the positive
$\mathbb Q[v,v^{-1}]$-basis existence remain valid.  N-004's conditional
completion theorem remains valid as an abstract statement, but its
integral freeness premise is now disproved for the actual rank-four
order.  The obstruction is mixed characteristic and is invisible to
a computation over $\mathbb Q(v)$ or to the characteristic-zero
specialized ambient image alone.

\subsection{N-001: the sign-point cell rotation does not extend to the full table}

The supplied counting computation gives a weighted cyclic symmetry on
the four labels in the inert cell at $D=0$.  In its notation the pairing
$\pi$ swaps $\gamma^2_1$ and $\gamma^2_2$ and fixes
$\gamma^2_{12}$ and $\gamma^2_{21}$.  The proposed identity is
\[
 c(x,y;z)w(z)=c(y,\pi z;\pi^{-1}x)w(\pi^{-1}x).
 \tag{N-1}
\]
Here is a decisive obstruction to extending that same pairing to a
permutation of all ten basis labels with every weight nonzero.

Set $e=\gamma^2_{12}$ and let $\mu$ be the unit.  The closed-form
multiplication table of the supplied counting consult, Theorem 2, gives
at $D=0$
\[
                 \mu e=e,\qquad e^2=16e.
 \tag{N-2}
\]
There is no other term in the latter product.  Any permutation extending
the cell pairing preserves that four-element set, so
$\pi^{-1}\mu$ is outside it.  Apply (N-1) to $(x,y;z)=(\mu,e;e)$.
The left side is $w(e)\ne0$, while the right side is zero because
$\pi e=e$ and $e^2$ has no coefficient at $\pi^{-1}\mu$.
This contradiction proves the claim for every possible extension and
choice of nonzero weights.  The certificate
\nolinkurl{local_obstructions_verified.json} verifies (N-2) directly.

This does not contradict cyclic invariance of the trilinear trace
$\tau(xyz)$: passing from that tensor to structure constants requires
inverting a Gram matrix, and the full form is degenerate here.  It also
does not deny the verified rotation on the cell alone.  It settles the
literal extension of that rotation to the full label table negatively.

\subsection{N-002: the first negative word has length five}

Every word in $P_1,P_2,P_3$ of length at most four lies in the
$C\otimes C$ cone.  The length-five word
\[
                         P_{12132}=P_1P_2P_1P_3P_2
\]
does not.  Its only coordinates with a negative Laurent coefficient are
\[
\begin{array}{c|c}
(a,b)&[C_a\otimes C_b]P_{12132}\\\hline
(2431,2314),(2314,2431)&2D(12-D)\\
(3421,2314),(2314,3421)&4\delta(8-D).
\end{array}
\tag{N-3}
\]
Their expansions are respectively
\[
 -2(v^4+v^{-4})+16(v^2+v^{-2})+36,
 \qquad -4(v^3+v^{-3})+20(v+v^{-1}).
\]
The exhaustive minimality check covers all $1+3+9+27+81=121$ words,
not just a sample.  The standard-$T$ replay script
\nolinkurl{verify_local_obstructions.py} constructs each word and then
performs triangular conversion to $C\otimes C$.  Thus the bounded
minimal-length assertion is exact.  The longer negative word of
INT-NSRH-4 is inherited and is not being recounted as new.

The sign mechanism can be seen directly.  Right dual-KL multiplication
has the descent coefficient $-\delta$; on an ascent its coefficients
are nonnegative KL $\mu$-coefficients.  Moreover
\[
 P_s=2C_s\otimes C_s+
       \delta(C_s\otimes1+1\otimes C_s)+D(1\otimes1).
 \tag{N-4}
\]
For the predecessor $P_{1213}$, the coordinates
$(2341,2134),(2341,2314),(2341,2341)$ are
$D(D+4),4\delta^3,8D$.  Their contributions to the first coordinate in
(N-3), upon right multiplication by $P_2$, are
$2D(D+4),-4D^2,16D$, whose sum is $2D(12-D)$.
The predecessor coordinates
$(3241,2314),(3241,2341),(3241,2134)$ are
$4D,8\delta,8\delta$; they contribute
$-4\delta D,16\delta,16\delta$ to the second coordinate.
The negative descent term has higher extreme Laurent degrees than the
positive remainder.  This explains the extreme-degree sign in this
word; it does not by itself prove a universal description of all
negative coordinates in F57 or F67.

Both polynomials in (N-3) change sign as real $v$ ranges over $(1,\infty)$.
Consequently neither has the form
$\pm(v-v^{-1})^a f(v)$ with $f\in\mathbb N[v,v^{-1}]$ nonzero and
$a\ge0$.  This excludes that factored positivity for this word in this
ambient basis, not for an unknown alternative basis of $\mathcal B_4$.

There is also an immediate categorical obstruction.  A split
Grothendieck realization whose effective indecomposable classes are
exactly the $C_a\otimes C_b$, whose grading shift is $v$, and in which
each actual $P_i$ is effective and composition is tensor multiplication,
cannot exist: effectiveness of their fivefold product contradicts
(N-3).  Derived or virtual realizations, different positive bases, and
dimension rings with cancelling pairs are outside that conclusion.

\subsection{N-003: an explicit positive generic basis and its exact scope}

An exact bounded search produces 114 bar-invariant elements of
$\mathcal B_4$ in the positive cone, forming a $K$-basis of
$K\mathcal B_4$.  Each is an integral combination of words of length
at most ten.  The complete recursive certificate is
\nolinkurl{mixed_cone_certificate.json}; it defines the elements without
large expanded polynomials as follows.  For its ordered list of words
$w_j$, set
\[
 B_j=P_{w_j}+\sum_{i<j}m_{ji}B_i,
 \qquad m_{ji}\in\{0,1,2,3,4,5,6\}.
 \tag{N-5}
\]
All 77 initial words are themselves positive.  The remaining 37 words
are repaired using the listed earlier positive elements. Individual additions have coefficients at most five; collecting repeated donors gives coefficients at most six. No power of
$D$ is needed in this successful certificate.  For example
\[
                    P_{12132}+P_{12323}+P_{121323}
\]
is positive.  Every correction donor has index at least 31 in the
zero-based certificate, so no correction uses the unit.

The bounded search first removes adjacent repetitions using
$P_i^2=DP_i$ and orders commuting $P_3P_1$ as $P_1P_3$.
The exact counts and cumulative ranks are
\[
\begin{array}{c|rrrrrrrrrr}
\text{length}&1&2&3&4&5&6&7&8&9&10\\\hline
\text{normal words}&3&5&8&13&21&34&55&89&144&233\\
\text{positive words}&3&5&8&13&17&18&20&21&24&30\\
\text{all-word rank}&4&9&17&30&49&75&101&113&114&114\\
\text{positive-word rank}&4&9&17&30&45&55&61&66&70&77
\end{array}
\tag{N-6}
\]
The ranks in (N-6) are certified lower bounds over $K$, computed by
row reduction at $v=2$ modulo 1000003; reaching the inherited upper
bound 114 proves equality.  The positive-coordinate assertions use
exact integer Laurent arithmetic, not numerical sign tests.

The search uses the direct dual-KL action (N-4).
\nolinkurl{verify_positive_basis.py} independently replays (N-5) using
standard-$T$ multiplication and only then converts to dual-KL
coordinates.  It checks every coefficient and modular independence of
all 114 elements.  Its output is
\nolinkurl{positive_basis_verified.json}.  Since $P_i$ and the two KL
bases are bar-invariant, (N-5) is bar-invariant as well.

Let $L=\sum_j A B_j$.  The triangular change (N-5) is unipotent, so
$L=\sum_j A P_{w_j}$ and these are two $A$-bases of $L$.
\emph{They are not bases of the whole order.}  At $v=2$, expand the
actual product $P_{232321}P_2$ in the 114 uncorrected words.  Its unique
coordinate at the basis word $P_1$ is
\[
                            -\frac{10656875}{4608}.
 \tag{N-7}
\]
The denominator is $2^9\cdot3^2$ and the numerator is not divisible by
3.  An $A$-linear combination would specialize to coefficients in
$\mathbb Z[1/2]$.  The specialized words are linearly independent, so
(N-7) is incompatible with such a combination.  Hence
$P_{232321}P_2\notin L$, and $L\subsetneq\mathcal B_4$.
The exact rational elimination and all coordinates are recorded in
\nolinkurl{lattice_specialization_certificate.json} and reproduced by
\nolinkurl{lattice_specialization.py}.  This is a proved obstruction to
this lattice, not a numerical suspicion or a proof against every
positive integral basis.

\subsection{N-004: the all-degree bare cone problem reduces to freeness}

The positive generic basis has a useful consequence independent of the
search degree.  Write $z=\sum_{j=0}^{113}B_j$.  Every ambient
$C_a\otimes C_b$ coordinate that occurs nontrivially anywhere in
$K\mathcal B_4$ has a nonzero coefficient in $z$, and that coefficient
is a nonnegative Laurent polynomial.  Indeed a coordinate vanishing
on every $B_j$ vanishes on their $K$-span; nonnegative coefficients
prevent cancellation in the sum.

There is an $A$-algebra augmentation
$\epsilon_-:\mathcal B_4\longrightarrow A$ sending every $P_i$ to
zero.  It is the restriction of the tensor product of the trivial
character on one standard Hecke factor and the sign character on the
other.  The certificate begins with the unit, all other words are
nonempty, and no correction uses the unit.  Therefore
\[
                              \epsilon_-(z)=1.
 \tag{N-8}
\]

\paragraph{The all-degree theorem.}
The following statements are equivalent:
(i) $\mathcal B_4$ is a free $A$-module;
(ii) $\mathcal B_4$ has an $A$-basis in the $C\otimes C$ cone.
Here ``basis in the cone'' imposes no cellularity, canonical normalization, bar normalization, requirement that the unit itself be a basis vector, or minimal degree requirement.

Only (i)$\Rightarrow$(ii) needs proof.  Under (i), $\mathcal B_4$ is
free of rank 114, and (N-8) splits it as
$Az\oplus\ker\epsilon_-$.  The kernel is stably free of rank 113.
A stably free module of rank greater than one over
$\mathbb Z[v,v^{-1}]$ is free; for example apply Abedelfatah (2010),
Theorem 1.3 and its stated consequence, with $R=\mathbb Z$.
Choose a basis $z,b_2,\ldots,b_{114}$.

For a fixed $b_j$, each negative coefficient occurs at one of finitely
many Laurent exponents in an active ambient coordinate.  For each
such coordinate choose one positive monomial of its coefficient in
$z$.  Choose $R$ large enough that every negative exponent is obtained
from its chosen positive exponent by a shift in $[-R,R]$.
Then, for a sufficiently large positive integer $N$,
\[
 b_j+N\left(\sum_{k=-R}^R v^k\right)z
 \tag{N-9}
\]
has no negative coefficient in any coordinate.  The added term has
nonnegative coefficients everywhere, so it creates no new negative
coefficient outside the chosen finite set.  Perform (N-9) independently
for every $j$.  These elementary column operations preserve an
$A$-basis, while $z$ is already positive.  This proves (ii) in every
degree, assuming precisely the freeness in (i).

\paragraph{An unconditional rational-coefficient corollary.}
The order $\mathbb Q\otimes\mathcal B_4$ is a finite torsion-free module
over the PID $\mathbb Q[v,v^{-1}]$, because it is a submodule of the
finite free ambient Hecke tensor algebra.  It is consequently free,
and the same argument gives a positive
$\mathbb Q[v,v^{-1}]$-basis.  Clearing constant rational denominators
in each resulting vector preserves a basis over that PID and puts all
vectors in $\mathcal B_4$ with nonnegative \emph{integer} Laurent
coordinates.  This existence theorem is not a certificate that the
particular 114 bounded-degree vectors above are a basis of this larger
rational order.

The integral freeness premise is disproved by N-009. This is a module obstruction on $\operatorname{Spec}\mathbb Z[v,v^{-1}]$: torsion-free over this two-dimensional ring does not imply locally free.  A generic basis, or a positive basis of a proper integral lattice, cannot remove that obstruction.
The GCT IV source explicitly proves freeness for rank three and gives
the generic dimension and a field-level monomial computation for rank
four; the latter is not an integral freeness proof.

\subsection{N-005: smaller composing dimension rings at rank three}

Set $q=v^2$ and $D=q+2+q^{-1}$.  We use the ten labels
$\mu,\Sigma,\gamma^i_x$, $i=1,2$, $x\in\{1,12,2,21\}$,
and the true trace constants of the supplied counting consult,
Theorem 2, with the corrected sign of $b_2$.
For each of the two choices $b=-8,-4$, put
\[
\begin{gathered}
 a=(D-1)(D-2),\quad W=7D+b,\quad
 \Lambda_1=D,\quad\Lambda_2=3D^2-10D+4,\\
 G_1=2Wa,\quad G_2=Wa,\quad S=2Wa^2,\\
 U_1=2W-D,\quad U_2=(W-D)/2,\\
 V_1=DU_1+D^2-D,
 \qquad V_2=DU_2+(D^2-\Lambda_2)/2.
\end{gathered}
\tag{N-10}
\]
All these are integer polynomials.  Give $\mu$ its unit products and set
\[
\begin{split}
 \Sigma^2&=S\Sigma,\\
 \Sigma\gamma^i_x=\gamma^i_x\Sigma&=D^{|x|-1}G_i\Sigma,\\
 \gamma^i_x\gamma^j_y&=D^{|x|+|y|-2}W\Sigma\quad(i\ne j).
\end{split}
\tag{N-11}
\]
For a same-cell product, reduce the word $xy$ using
$P_i^2=DP_i$ and $P_1P_2P_1=\Lambda_iP_1$,
$P_2P_1P_2=\Lambda_iP_2$, to
$D^\alpha\Lambda_i^\beta P_z$, with $\beta\in\{0,1\}$.
Define
\[
 \gamma^i_x\gamma^i_y
  =aD^\alpha\Lambda_i^\beta\gamma^i_z+
 \begin{cases}
 D^{|z|+\alpha-2}U_i\Sigma,&\beta=0,\\
 D^{|z|+\alpha-1}V_i\Sigma,&\beta=1.
 \end{cases}
 \tag{N-12}
\]
Every exponent of $D$ in (N-12) is nonnegative for the listed four
words.  These formulas give 132 nonzero structure constants,
including the 19 unit constants.

\paragraph{Associativity.}
Work first over $\mathbb Q(D)$.  For each $\Lambda_i$, the four-dimensional
cell algebra on the four alternating words is associative: for example
realize its two generators as
\[
 \begin{pmatrix}D&1\\0&0\end{pmatrix},\qquad
 \begin{pmatrix}0&0\\\Lambda_i&D\end{pmatrix}.
\]
Their four nonempty alternating words of lengths one and two are
linearly independent over $\mathbb Q(D)$ for the parameters in (N-10),
and multiplication has exactly the indicated reductions.  Scaling
these basis vectors by $a$ gives the cell terms in (N-12).
Use the ten-dimensional product of the two matrix algebras and two scalar algebras, the latter called the character and augmentation summands. Send $\Sigma$ to
$S$ in the scalar summand and to zero in the two cell summands and the augmentation summand; send
$\gamma^i_x$ to its scaled cell vector and to
$D^{|x|-1}G_i$ in the character summand and zero in the augmentation summand; send $\mu$ to the unit in all four factors. This map is injective: the augmentation detects $\mu$, then the two matrix factors detect the eight cell vectors, and the character factor detects $\Sigma$.

The necessary scalar residual in a same-cell product is
\[
 \frac{D^{|x|+|y|-2}G_i^2-
 aD^\alpha\Lambda_i^\beta D^{|z|-1}G_i}{S}.
\]
Since $|x|+|y|=|z|+\alpha+2\beta$, this is exactly (N-12), using
$G_i(G_i-aD)/S=U_i$ and
$G_i(DG_i-a\Lambda_i)/S=V_i$.
The cross-cell residual is $G_1G_2/S=W$.
Thus the table embeds in a product of associative algebras and is
associative over $\mathbb Q(D)$.  Its entries lie in $\mathbb Z[D]$,
so every associativity identity already holds there.

\paragraph{Effectivity, domination and parity.}
Write $c$ for the true table and $d_b$ for (N-10)--(N-12).
The true cell parameters are
$\delta_1=-(D-1)(D-4)$, $\delta_2=(D-1)(D-4)$,
$\lambda_1=D$, $\lambda_2=(D-2)^2$.
Coefficientwise in $q$,
\[
 a=|\delta_i|+2,\qquad
 \Lambda_2-\lambda_2=2(q^2+q+q^{-1}+q^{-2}).
\]
Thus the cell products dominate the absolute trace coefficients and
have the same parity.  Products of nonnegative polynomials preserve
these inequalities, since $|fg|\le |f|\,|g|$ coefficientwise.
For the bottom and scalar products it suffices to check the following
half-differences.  A displayed list is centered at exponent zero.
\[
\begin{array}{c|c|c}
\text{half-difference}&b=-8&b=-4\\\hline
(W-D)/2&(3,2,3)&(3,4,3)\\
(U_1-D)/2&(6,4,6)&(6,8,6)\\
(U_2-D)/2&(1,0,1)&(1,1,1)\\
(V_1-(5D-4))/2&(7,17,22,17,7)&(7,21,30,21,7)\\
(V_2-(2D^2-5D+4))/2&(3,3,3)&(4,5,4)\\
(G_1-4D(D-1))/2&(7,11,21,18,21,11,7)&(7,15,25,26,25,15,7)\\
(G_2-D^2(D-1))/2&(3,4,8,6,8,4,3)&(3,6,10,10,10,6,3)
\end{array}
\tag{N-13}
\]
The remaining half-difference $(S-4D^2(D-1)^2)/2$ has lists
\[
\begin{array}{c|l}
-8&(7,18,42,52,67,60,67,52,42,18,7),\\
-4&(7,22,50,72,91,92,91,72,50,22,7).
\end{array}
\]
These are nonnegative integers.  This proves all of
$d_b\ge |c|$, $d_b\equiv c\pmod2$, and nonnegativity for the full table.
For $b=-4$, every basic polynomial used multiplicatively in the table
is symmetric and unimodal, except that $\Lambda_2$ itself need not be;
it only occurs multiplied by $a$, and
\[
 a\Lambda_2=(3,5,10,11,14,11,10,5,3)
\]
is symmetric unimodal.  Products of nonnegative symmetric unimodal
Laurent polynomials are symmetric unimodal, by decomposing each into
nonnegative sums of centered interval polynomials.  Thus all 132
constants for $b=-4$ are unimodal.  The $b=-8$ table has 24 failures of
unimodality, recorded in the certificate.

\paragraph{The exact cancelling-pair scores.}
The score is $\frac12\sum_{x,y,z,k}(d_{xyz,k}-|c_{xyz,k}|)$.
The unit constants have zero excess.  The sum of absolute coefficients
of all nonunit true constants is 4424: the bottom/scalar terms
contribute 4096 and the two cell terms contribute 152 and 176.
For clarity, the same-cell reduction types $(\alpha,\beta)$ have
multiplicities $(0,0):2$, $(1,0):6$, $(0,1):6$, $(1,1):2$.
The corresponding absolute norms in cell one are $4,8,8,24$ and in
cell two $4,8,16,12$, giving those totals.
At $q=1$, hence $D=4$, the sum of nonunit dimension coefficients is
\[
 S+20(G_1+G_2)+200W+20(U_1+U_2+V_1+V_2)
   +2a(7\Lambda_1+13)+2a(7\Lambda_2+13).
 \tag{N-14}
\]
It is 18976 for $b=-8$ and 22504 for $b=-4$.  Consequently the scores are
\[
 \boxed{7276\quad(b=-8)},\qquad
 \boxed{9040\quad(b=-4,\ \text{all dimensions unimodal})}.
 \tag{N-15}
\]
The corresponding bottom-cell scores are 6612 and 8376.
They are strictly below the inherited window minima 7952 and 11764.
Thus neither earlier minimum is global.  Nothing here proves that
7276 or 9040 is globally minimal or that minimization selects a
canonical ring.

The independent supplied-table replay
\nolinkurl{verify_dimension_improvements.py} first checks the true table
against all 132 stored constants, then checks both new tables on all
1000 ordered triples and every Laurent coefficient.  Its detailed
output is \nolinkurl{dimension_improvements_verified.json}.

\paragraph{Effective classes and the categorical qualification.}
In the coefficient ring
$\mathbb Z[q,q^{-1},\varepsilon]/(\varepsilon^2-1)$, each pair
$(d_b,c)$ lifts to the effective class
\[
                 \widehat c=(d_b+c)/2+\varepsilon(d_b-c)/2.
\]
Evaluation at $\varepsilon=1,-1$ embeds that ring into pairs of
Laurent polynomials of equal parity, so associativity of both tables
proves associativity of these effective lifts.  The left regular
matrix realization in the supplied counting consult, Theorem 5,
then gives its specified abstract monoidal realization.  The objects
can be interpreted as graded sign Artin--Tate objects.  No equivalence
with the Hecke category, canonical choice of basis, or canonical
choice of $d_b$ has been established by this construction.

\subsection{Remaining geometric and canonical-basis gates}

\paragraph{N-006: a literal fixed-parameter sheaf and its limitation.}
Fix $q>1$ and evaluate the actual End/Hom model of GCT\_NSRH F46 at
$v=q^{1/2}$.  Its specified finite Weil spaces $H,K$ have dimensions
6 and 44; $\operatorname{End}_F(H)$ is the actual rank-three algebra
and $\operatorname{Hom}_F(H,K)$ is its actual 64-dimensional right
module, with the named $P_1,P_2$ actions.  Its eigenvalues are
$-1,i,1,\lambda$, where
\[
 \lambda=(1+i\Delta)/(1-i\Delta),\qquad
 \Delta=(D-1)(D-4)/D^2\in[0,1).
\]
Every complex conjugate has modulus one.  Choose an admissible
$\ell\ne\operatorname{char}\mathbb F_q$ and embeddings of these
algebraic coefficients in $\overline{\mathbb Q}_\ell$; for a genuine
lisse sheaf one may choose $\ell$ outside the finite set where an
eigenvalue is not a unit.  A finite Weil sheaf requires only the
specified invertible Frobenius.

For any of these pure weight-zero coefficient representations $C$,
let $\pi:\mathbb A^1_{\mathbb F_q}\to\operatorname{Spec}\mathbb F_q$
and put $\mathscr F_C=\pi^*C(1)$.  At a closed point of degree $d$,
its eigenvalues are $q^{-d}\alpha^d$, so its stalk is pure of weight
$-2$, checked directly without using compact cohomology.
The standard affine-line calculation gives
\[
 H_c^2(\mathbb A^1_{\overline{\mathbb F}_q},\mathscr F_C)=C,
 \qquad H_c^j=0\quad(j\ne2).
\]
All specified endomorphisms lift as constant sheaf morphisms, hence
retain the actual $P_i$ actions.  More explicitly, pullback from the
arithmetic point is fully faithful on these constant arithmetic local
systems: evaluation at a rational point recovers the intertwiner,
and a morphism is exactly a linear map commuting with Frobenius.
Consequently
\[
 \operatorname{End}(\mathscr F_H)=A,\qquad
 \operatorname{Hom}(\mathscr F_H,\mathscr F_K)=M
\]
with the exact named algebra and precomposition action of F19/F21/F46.
The common Tate twists cancel in these morphism spaces.
The F46 intertwiner space $M$ has identity Frobenius, and
$\mathscr F_M=\pi^*M(1)$ directly satisfies
$H_c^2(\mathbb A^1,\mathscr F_M)=M$, with the actual $P_1,P_2$
action.  This is the specified tensor module, not a substitute of
matching dimension: F21 supplies the regular action-preserving
basis comparison, and F46 explicitly retains its generator matrices.

Thus the literal fixed-$q$ sheaf request in item 8 is solved for the
actual F19/F21 algebra and tensor action and the F46 coefficient
objects.  The proof is a formal affine-line construction from the
explicit F46 spectra.  F44 also realizes this same algebra and
module as a coherent corner; its separately specified semilinear
ring Frobenius is not identified with F46's chosen linear Weil
operator by this construction.

If the request additionally requires the same finite coefficient
space and operators while replacing $v$ by $q^{d/2}$ at every extension,
the inherited F26/F42 obstruction applies.  The fixed relation
$P^2=D(q)P$, together with the demanded retuned nonzero traces at
powers one and two, would force $D(q^d)=D(q)$, false for $d=2$.
Thus that stronger fixed-object interpretation is impossible.
Extension-wise retuning is an additional property, not a condition of
the literal fixed-$q$ request just solved.  A varying geometric family
compatible with that retuning is a separate question.

\paragraph{N-007: rational descent does not ensure a finite Tate polynomial.}
The supplied ground-ring theorem gives rational functions in $D$ for
a $\mathbb Q(D)$-basis of words.  Rationality alone does not imply
Laurent-polynomial regularity.  Already in rank two, the basis
$\{1,b\}$ with $b=P_1/D^2$ has
\[
                        b^2=(1/D)b,
 \qquad 1/D=q/(q+1)^2.
\]
This coefficient is bar-invariant and rational in $D$, but has a pole
at $q=-1$ and is not a Laurent polynomial.  It cannot be a finite
integer-weight Tate Euler polynomial.  Completed or rational
stack-theoretic Tate expressions are a different notion.  The
rational ground-ring statement remains valid; the finite-polynomial
conclusion needs additional normalization or regularity.

\paragraph{N-008: the sign-point ambient image is not the abstract fiber.}
Exact closure and rational row reduction at $v=i$ give dimension 55
for the image of the rank-four $P$-generated algebra in the specialized
ambient tensor Hecke algebra.  Its restricted standard trace Gram
matrix has rank 42.  The same calculation at rank three gives image
dimension 7 and Gram rank 6, whereas the abstract rank-three order
has ten-dimensional fiber.  These are image calculations, not a
rank-four abstract fiber computation or an $E(4)$ multiplication table.
The script \nolinkurl{sign_specialization.py} uses exact integer
coordinates after absorbing the known powers of $i$, closes the span
under every generator, and records the full basis and Gram matrix.
No claim about a Jacobson radical follows from the Gram rank alone.

\paragraph{Why $E(4)$ is still a gate.}
The supplied GCT VIII, sections 3.1 and 3.3, specifies an inductive
nonstandard Gelfand--Tsetlin framework, type order, and conjectural
minimum-degree/uniqueness conditions.  Its section 5.2 reports the
rank-four distinguished element $\Sigma_4$, not a certified full
114-element table.  GCT VII does not provide the missing explicit
rank-four Jucys--Murphy operator formula alleged in an earlier intake.
Blasiak (2012), Theorem 5.1 and section 6, gives irreducibles and
seminormal bases for the nonstandard Temperley--Lieb quotient; that
quotient is not the full 114-dimensional algebra.
The present bounded cone basis is not identified with VIII's $E(4)$.
Consequently its field of definition, all signed structure constants,
$\mathbb N[q+q^{-1}]$ property, and a dimension ring for that designated
basis have not been determined here.  Rational descent for arbitrary
word spans cannot settle those normalization questions.

\paragraph{Source locators and novelty limits.}
GCT IV: supplied \texttt{Apr13KroneckerGCT4.tex}, Remark 11.19 and
subsection \texttt{sb4}; the former explicitly calls rank-four cone
positivity uncertain.  The closed rank-three table and dimension-ring
criterion are inherited from
\nolinkurl{OURS_CONSULT_a157422a_COUNTING_CONSTRUCTION.md},
Theorems 2, 5, 7, 7$'$ and section 4.1.  The cell-only sign symmetry is
in \nolinkurl{OURS_E4_RANK_FOUR_20261121.md}, section 6b.
The pure fixed-parameter model and all-extension obstruction are
GCT\_NSRH F46 and F26/F42.  The new finite certificates and theorems
above were not located in those supplied sources; this is not a claim
of a comprehensive worldwide novelty proof.

For the cancellation theorem used in N-004 see Abedelfatah, A. (2010).
\emph{On stably free modules over Laurent polynomial rings},
arXiv:1012.3540v1, Theorem 1.3 and abstract, p. 1,
\texttt{https://arxiv.org/abs/1012.3540}.
For the quotient distinction see Blasiak, J. (2012).
\emph{Representation theory of the nonstandard Hecke algebra},
arXiv:1201.2209v2, Theorem 5.1 and section 6.
For affine-line cohomology use Milne, J. S. (2013),
\emph{Lectures on etale cohomology}, version 2.21, Proposition 18.3
and Theorem 24.1; the stalk weight assertion here is checked directly
from the displayed eigenvalues.


\subsection{N-010: an explicit split form and GT path basis over $\mathbb Q(D)$}
\label{sec:nonstandard-rational-gt}

\paragraph{Statement and scope.}
Let $F=\mathbb Q(D)$, $D=(v+v^{-1})^2$, and let
$\mathcal B_{r,F}$ be the $F$-span of the actual $P$-words in the
generic nonstandard algebra.  The fact that this is an $F$-form is the
supplied counting consult, Theorem 1.  The full algebra at rank four is
split over $F$:
\[
 \mathcal B_{4,F}\cong
 F\times M_2(F)\times M_3(F)\times M_6(F)
 \times M_5(F)\times M_5(F)\times M_2(F)\times M_3(F)\times F.
 \tag{N-25}
\]
The tower through rank four has multiplicity-free branching, and every
GT path line, together with a concrete choice of vector on it, can be
constructed over $F$.  The accompanying certificate gives all nine
triples of generator matrices, all GT transition matrices, all path
labels, and an exact rank certificate for (N-25).

This does not identify these choices with VIII's $E(4)$.  It proves that
neither splitting the algebra nor choosing its GT path vectors forces
an algebraic extension of $\mathbb Q(D)$ at rank four.  The separate
minimum-degree globalization and its prescribed normalization remain
uncomputed.

\paragraph{Three small Hecke models.}
Put $\delta=v+v^{-1}$.  The following matrices represent
$A_i=\delta C'_{s_i}$, with column-vector convention:
\[
 A_i^{(4)}=(D)\qquad(i=1,2,3),
\]
\[
 A_1^{(31)}=\begin{pmatrix}0&0&0\\D&D&0\\0&0&D\end{pmatrix},\quad
 A_2^{(31)}=\begin{pmatrix}D&1&0\\0&0&0\\0&1&D\end{pmatrix},\quad
 A_3^{(31)}=\begin{pmatrix}D&0&0\\0&D&D\\0&0&0\end{pmatrix},
 \tag{N-26}
\]
\[
 A_1^{(22)}=A_3^{(22)}=
 \begin{pmatrix}D&1\\0&0\end{pmatrix},\qquad
 A_2^{(22)}=\begin{pmatrix}0&0\\D&D\end{pmatrix}.
 \tag{N-27}
\]
For each triple, direct multiplication verifies
\[
 A_i^2=DA_i,\qquad A_1A_3=A_3A_1,\qquad
 A_iA_{i+1}A_i-DA_i=A_{i+1}A_iA_{i+1}-DA_{i+1}.
 \tag{N-28}
\]
Consequently $T_i=\delta^{-1}A_i-v^{-1}I$ obey the ordinary
Hecke quadratic and braid relations, so these are actual Hecke
representations after extension to $\mathbb Q(v)$.  In the braid
calculation, putting $a=\delta v^{-1}$ gives
$a(a-D)=-D$, which proves the equivalence with the last relation
in (N-28).  The quadratic relation follows from $A_i^2=DA_i$.

For provenance, the script reconstructs (N-26)--(N-27) from the supplied
exact KL arithmetic.  On a $C'$-cell basis indexed by permutations $w$,
rescale by $\delta^{\ell(w)\bmod2}$.  A diagonal coefficient $\delta$
in $C'_i$ becomes $D$ in $A_i$; an off-diagonal integer coefficient
joining opposite parities becomes an integer or an integer times $D$.
The script checks these assertions coefficient by coefficient.  In this
$C'$ convention the RSK cell shape is the transpose of the displayed
Specht label: the identity cell is the sign representation.  More
fundamentally, the explicit checks (N-28) suffice for the construction,
independently of a general cell-module theorem.

\paragraph{Nine actual nonstandard modules over $F$.}
For two triples $A=(A_i)$ and $B=(B_i)$ above, the actual tensor
generator $P_i=C'_i\otimes C'_i+C_i\otimes C_i$ acts as
\[
 \mathsf P_i(A,B)=\frac{2}{D}A_i\otimes B_i
              -A_i\otimes I-I\otimes B_i+D I.
 \tag{N-29}
\]
Every entry is in $F$.  For $A=B$, (N-29) commutes with the tensor flip;
hence its symmetric and alternating subspaces are defined over $F$.
The symmetric square contains a one-dimensional simultaneous
$P_i=D$ eigenspace.  The script computes that line exactly over $F$
and takes the quotient, denoted $S'M_{\lambda}$; its dimensions are
two for $\lambda=(2,2)$ and five for $\lambda=(3,1)$.

Twisting by $\Theta(P_i)=D-P_i$ is also defined over $F$.
Indeed, the ordinary Hecke involution $T_i\mapsto-T_i^{-1}$ on one
tensor factor sends $C'_i\mapsto-C_i$ and $C_i\mapsto-C'_i$, so it
sends $P_i$ to $D-P_i$.  This proves validity of the twist in the
full ambient algebra, not merely in a quadratic presentation.

The nine constructed modules, in the order of (N-25), are
\[
 \epsilon_-,\ S'M_{22},\ \Lambda^2M_{31},\ M_{31}\otimes M_{22},\
 S'M_{31},\ (S'M_{31})^\Theta,\ M_4\otimes M_{22},\
 M_4\otimes M_{31},\ \epsilon_+.
 \tag{N-30}
\]
Here $\epsilon_-$ has $P_i=0$, and $\epsilon_+$ has $P_i=D$.
These names agree with the known list in GCT IV, subsection 11.8,
printed p.~67.  The construction itself provides a homomorphism
from $\mathcal B_{4,F}$ to the product of the nine endomorphism
algebras.  There is no circular assumption that these modules are
simple or pairwise distinct: that follows from the rank check next.

\paragraph{Exact splitting certificate.}
Flatten the nine matrix blocks of an element, in their displayed order,
to a vector with $1+4+9+36+25+25+4+9+1=114$ coordinates.  Specialize
$D=7$ and reduce modulo the prime $1000003$.  All denominators occurring
in the generator matrices are invertible there.  Exact Gaussian
elimination of actual word vectors gives rank $114$; the certificate
lists the 114 accepted words and every pivot column.  Its longest word
has length nine.  A determinant that is nonzero after a valid
specialization is nonzero in $F$, so the generic image has dimension
114.  The inherited generic dimension of the source is also 114.
The homomorphism is therefore an isomorphism, proving (N-25), absolute
simplicity of each factor module, and pairwise nonisomorphism.
This is a finite rank proof of the complete algebra decomposition,
not an inference that a short word search presents the integral order.

\paragraph{The exact branching projectors.}
Write $N_1,N_2$ for the two two-dimensional rank-three simples,
with $P_1P_2P_1=\lambda_iP_1$, where
$\lambda_1=D$ and $\lambda_2=(D-2)^2$.  In each of the nine models put
\[
 Z=(P_1-P_2)^2,\qquad z_0=0,\quad z_1=D(D-1),\quad z_2=4(D-1).
\]
The supplied exact arithmetic checks
$Z(Z-z_1I)(Z-z_2I)=0$ and that $Z$ commutes with $P_1,P_2$.
The distinct roots give the spectral idempotents
\[
 E_j=\prod_{k\ne j}\frac{Z-z_kI}{z_j-z_k}\quad(0\le j\le2).
 \tag{N-31}
\]
The exact checks also verify $E_0(P_1-P_2)=0$ and
$E_j(P_1P_2P_1-(D^2-z_j)P_1)=0$ for $j=1,2$.
On the zero eigenspace, $E_0P_1/D$ and $E_0(I-P_1/D)$ are the
$\epsilon_+$ and $\epsilon_-$ projectors.  The other two are the
$N_1$ and $N_2$ projectors.  Their exact ranks yield:
\[
\begin{array}{c|l}
\mathcal B_{4,F}\text{-module}&\text{restriction to }\mathcal B_{3,F}\\\hline
\epsilon_-&\epsilon_-\\
S'M_{22}&N_2\\
\Lambda^2M_{31}&\epsilon_-\oplus N_1\\
M_{31}\otimes M_{22}&\epsilon_+\oplus\epsilon_-\oplus N_1\oplus N_2\\
S'M_{31}&\epsilon_+\oplus N_1\oplus N_2\\
(S'M_{31})^\Theta&\epsilon_-\oplus N_1\oplus N_2\\
M_4\otimes M_{22}&N_1\\
M_4\otimes M_{31}&\epsilon_+\oplus N_1\\
\epsilon_+&\epsilon_+
\end{array}
 \tag{N-32}
\]
Each constituent occurs once.  For rank three to rank two, each $N_i$
has one $P_1=D$ and one $P_1=0$ line.  Therefore the nonzero projectors
$E_jP_1/D$ and $E_j(I-P_1/D)$, together with the one-dimensional
constituents already identified, give every GT path line over $F$.

The certificate chooses the first nonzero column of each rank-one
projector and divides it by its first nonzero entry.  The resulting
transition matrix is invertible over $F$.  The script checks that all
conjugated $P_i$ matrices agree exactly and that $P_1$ is diagonal with
the specified entries $D$ or $0$.  This fixes a reproducible rational
normalization and path order for this construction.  Under (N-25),
the associated matrix units are $F$-linear combinations of the
certified word basis, by solving the invertible 114-dimensional
linear system over $F$.  Thus they give a full GT matrix-unit basis
inside $\mathcal B_{4,F}$.

\paragraph{The remaining canonical-basis gate, precisely.}
GCT VIII, subsection 3.1, printed pp.~19--20, describes the GT path
construction under its multiplicity-free condition C1'.  Formula
(N-32) verifies that condition through rank four for the full algebra.
But VIII, subsection 3.3, printed pp.~20--22, additionally requires
the order induced from nonstandard duality, the scaling exponents
$a_T$, and unique global minimum-degree lifts.  These are its
Conjectures 3.1 and 3.2, not consequences of the split algebra.
Its subsection 5.2 computes a single distinguished rank-four element;
it supplies no complete rank-four table.

There is also a basic normalization issue that the degree comparison
alone cannot remove.  Suppose an initial GT vector $c'_T$ is replaced
by $u_Tc'_T$, with $u_T\in\mathbb Q^\times$.  Such changes arise, for
example, from a rational diagonal change of basis in an irreducible
module, with the dual factors transformed inversely.  Every lattice
generated in VIII's construction remains the same, because $u_T$ is
a unit in all the coefficient rings.  The residue label changes to
$u_Tb_T$, and its fibre is the scalar multiple by $u_T$ of the old
fibre.  For every candidate $y$, the ambient degree tuple satisfies
$d(u_Ty)=d(y)$.  Consequently, if the unique minimizing lift exists,
it changes to $u_Te_T$; the optimization in $a_T$ has the same
equivariance.  Thus a uniquely normalized table needs specified
initial constant normalizations or a further normalization rule.
The degree comparison by itself does not select them.

The present matrices establish neither that VIII's eventual $E(4)$
belongs to $\mathbb Q(D)$ nor that its signed factors lie in
$\mathbb N[q+q^{-1}]$.  They do show that an algebraic coefficient
extension cannot be justified merely by the need to split the
rank-four algebra or its GT tower.  Matrix units give structure
constants zero and one for their own $F$-basis, but this is no
dimension-ring assertion for VIII's designated basis and no integral
$A$-basis assertion, which N-009 excludes.

\paragraph{Literature status and verification.}
The nine generic simple models and generic dimension are known from
GCT IV.  GCT VII, subsection 7.2 (Example 2), explicitly states
splitting over its $\mathbb Q(q)$ in the paragraph following the
74-term relation figures, printed p.~40; its explicit simple-module
formulas follow in subsection 7.2.1.  Blasiak's \emph{Representation
theory of the nonstandard Hecke algebra}, Theorem 5.1 and section 6,
prove multiplicity-free branching and seminormal-basis results for
the nonstandard Temperley--Lieb quotient.  That quotient has generic
rank 89 at rank four, and it must not replace the full rank-114 algebra.
The contribution here is the explicit $\mathbb Q(D)$ descent, exact
full-algebra rank certificate, and complete rational GT path data;
the classification itself is not claimed as new.

The file \nolinkurl{rational_gt_rank4.py} contains the entire exact
construction and writes \nolinkurl{rational_gt_rank4_verified.json}.
It checks (N-28), symmetric/alternating restriction, quotient invariance,
the finite-field rank, all projectors, all path ranks, and all GT
intertwining identities.  The plethysm branch independently read the complete proof and script,
replayed the full construction, and checked the additional scalar and
rank-three parameter identities above.  Its whole-proof review is
\nolinkurl{nonstandard_plethysm_review_N010.md}, with the separate
additional-check receipt
\nolinkurl{rational_gt_rank4_plethysm_additional.json}.  The two
inherited inputs (generic dimension and $F$-form) retain that status.


\subsection{N-011: why the named obstructions occur at the Laurent edges}
\label{sec:nonstandard-edges}

The inherited negative-coordinate and trace obstructions are not being
claimed again.  The following elementary transfer identity explains
their locations and distinguishes the two different edge patterns.
Let
\[
 f(D)=\sum_{j=0}^m a_jD^j,\quad a_m\ne0,\qquad
 g(v)=\delta^\varepsilon f(D),\quad
 \delta=v+v^{-1},\quad D=\delta^2,\quad\varepsilon\in\{0,1\}.
\]
Put $L=2m+\varepsilon$.  Since
$\delta^n=v^n(1+v^{-2})^n$, the coefficient of $v^{L-2r}$ is
\[
 [v^{L-2r}]g(v)=
 \sum_{k=0}^{\min(r,m)}
 a_{m-k}\binom{L-2k}{r-k},
 \tag{N-33}
\]
where an out-of-range binomial coefficient is zero.  In particular
the outermost and next coefficients are
\[
             a_m,\qquad La_m+a_{m-1}.
 \tag{N-34}
\]
The coefficients at the other end are equal by bar invariance.
Thus the highest $D$-term alone controls the extreme Laurent degree;
lower $D$-degrees cannot repair its sign.  At the next Laurent degree,
the second-highest $D$-term competes with the binomial contribution
from the highest one.

For the exact examples requested in item 7, this gives:
\[
\begin{array}{l|r|r|r|r|r}
\text{example}&L&a_m&a_{m-1}&[v^L]&[v^{L-2}]\\\hline
\text{GCT\_NSRH F57}&18&-1&20&-1&2\\
\text{INT-NSRH-4: }(2013,3201)&15&-5&102&-5&27\\
\text{INT-NSRH-4: }(2013,3021)&16&-3&84&-3&36\\
\text{INT-NSRH-4: }(2103,3021)&15&-3&70&-3&25\\
\text{F67, three-plus-three}&12&2&-36&2&-12\\
\text{F67, one-plus-three}&8&2&-20&2&-4
\end{array}
 \tag{N-35}
\]
The three INT-NSRH-4 pairs use that intake's zero-based permutation
notation; transposing each pair supplies its other three coordinates.
The exact full $D$-polynomials and Laurent lists are in
\nolinkurl{edge_coefficients_verified.json}.  For example the first is
\[
\delta\bigl(-5D^7+102D^6-515D^5+1856D^4-3888D^3
                     +5568D^2-4608D+2048\bigr).
\]
All three have a negative highest $D$-coefficient, followed by a
sufficiently large positive next coefficient.  Their negative Laurent
coefficients are consequently at the outermost positions; the exact
remaining coefficient lists certify positivity inward from there.
For F57 the inherited polynomial starts $-D^9+20D^8-\cdots$, giving
the same mechanism.  This is compatible with the direct mixed-descent
calculation in N-002: virtual negative terms can dominate the highest
degree while positive lower-degree terms contribute only farther in.

F67 is a different pattern: its outermost coefficients are positive,
and its negative coefficients occur one step inward.  For the
three-plus-three example, the inherited direct four-dimensional
tensor calculation gives
\[
 \tau_6(D)=D^6+(-D^3+18D^2-48D+32)^2
          =2D^6-36D^5+\cdots.
\]
The negative cross term in the square outweighs the binomial spill
from $2D^6$: $-36+12\cdot2=-12$.  For the one-plus-three supplement,
\[
 \tau_{1,3}(D)=2D^4-20D^3+84D^2-128D+64,
\]
and the analogous coefficient is $-20+8\cdot2=-4$.
The source tensor rule includes the complementary, virtual generators
$mI-N_i$.  Its whole-product trace has those negative cross terms
despite the input matrices being effective.  A positive value of the
sum-of-squares trace at a real parameter does not imply coefficientwise
positivity.

These explanations concern the named inherited examples.  They prove
no theorem that every negative coordinate or word trace must have its
negative coefficients at the edges.  In particular, F67 does not
support a claim that the outermost coefficient itself must be negative.

\paragraph{Provenance and exact checks.}
The inputs are GCT\_NSRH F57; INT-NSRH-4 in
\nolinkurl{INTAKE_GCT_NSRH_20261121_part3.md}, sections 1--2;
and GCT\_NSRH F67 with its one-plus-three supplement.  The script
\nolinkurl{edge_coefficients.py} recomputes the three ambient coordinates
to obtain their new $D$-expressions, verifies that they and their
transposes are precisely the six bad coordinates of the given word,
and checks (N-33) at every exponent of all six displayed cases.
For F67 it also derives the scalar polynomials directly from the four
by four output matrices.  The representation-theoretic occurrence
claims and original obstructions retain their inherited finding IDs.


\subsection{N-012: the printed variable-scaling minimum has no minimum}
\label{sec:nonstandard-viii-scaling}

\paragraph{Precise source statement.}
GCT VIII, subsection 3.3, Conjecture 3.2, printed p.~22
(arXiv:0709.0751v2), varies the rational exponent $a$ in
$c_{a,T}=q^a c'_T$ and minimizes ambient degree over the fibre
\[
 \psi_{a,T}^{-1}(b_{a,T})\subset
 R^{a,\le T}\cap\bar R^{a,\le T}\cap\mathcal B^{\le T}_{\mathbb Q}.
 \tag{N-36}
\]
Here $b_{a,T}$ is the class of $c_{a,T}$ modulo
$q^{1/2}R^{a,\le T}$, and the degree of an ambient coefficient is
its pole order at $q=\infty$ (subsection 3.3, equation (13), p.~21).
The displayed prescription ranges over all rational $a$ and does
not intersect (N-36) with the bar-fixed subspace.  Its surrounding
prose states bar invariance as a goal of globalization; that goal is
not the additional equation $y=\bar y$ in the displayed domain.

Under these literal definitions, the variable-scaling minimum does
not exist.  The defect already occurs on a one-dimensional algebra,
and it occurs on any minimal GT line of the full rank-four algebra.
This is a failure of the displayed prescription.  It does not assert
that no canonical generic basis can be defined by a corrected rule.

\paragraph{The rank-one counterexample.}
Take $\mathcal B=\mathbb Q(q^{1/2})$, $c'_T=1$ and its Laurent
form $\mathbb Q[q^{1/2},q^{-1/2}]$.  Write $A_0,A_\infty$ for
the rings regular at zero and infinity.  For any integer $n\ge0$ set
\[
 a_n=-n,\qquad y_n=q^{-n},\qquad
 R^{a_n}=q^{-n}A_0,\qquad\bar R^{a_n}=q^n A_\infty.
 \tag{N-37}
\]
The element $y_n$ lies in the Laurent form and in $R^{a_n}$;
its quotient by the generator of $R^{a_n}$ is exactly one, so
it maps to $b_{a_n,T}$.  Its quotient by the generator of
$\bar R^{a_n}$ is $q^{-2n}$, which is regular at infinity.
Thus $y_n$ is an admissible candidate in (N-36).  Its unique
nonzero ambient coefficient has degree $-n$.  A nonzero algebraic
function has a finite pole order, so no candidate can have degree
at most $-n$ for every $n$.  Conjecture 3.2 is therefore false
under its displayed unrestricted domain.

This is not an objection based on a scalar-normalization ambiguity:
the initial vector is fixed to be exactly $1$, and the exponents
are ordinary integers already contained in the stated coefficient
field.  Nor does an integral Laurent form remove the example;
every $q^{-n}$ belongs to that form.

\paragraph{The same defect at rank four, for every label order.}
Use the explicit rational GT path lines of N-010, and write
$F=\mathbb Q(D)$, $K=\mathbb Q(v)$ with $q=v^2$.
Let $T$ be any minimal label in the finite partial order used for
globalization.  Then
\[
                  \mathcal B^{\le T}=Kc'_T
 \tag{N-38}
\]
is one-dimensional.  The line is defined over $F$ by N-010, so
it has a nonzero bar-fixed generator $d_T\in\mathcal B_{4,F}$.
Indeed every $P$-word is bar invariant and every scalar in $F$
is fixed by bar.  Clearing the finitely many denominators in an
$F$-word expression gives
\[
 t=g(D)d_T\in\mathbb Q[D]\langle P_1,P_2,P_3\rangle
                  \subseteq\mathcal B_{4,\mathbb Q},
 \qquad \bar t=t,\qquad t\ne0,
 \tag{N-39}
\]
for some nonzero $g(D)\in\mathbb Q[D]$.
The last inclusion uses precisely VIII's ring generated by the
rescaled $P_i$ over its Laurent coefficient form; no saturation
or localization by $g(D)$ is used.

First choose the initial GT vector $c'_T=d_T$ and make $g$ monic,
of degree $m$.  Put
\[
                 a_n=-m-n,\qquad y_n=q^{-n}t
                      \quad(n\ge0).
 \tag{N-40}
\]
Since $D=q+2+q^{-1}$,
$q^m g(D)$ is regular at zero and has value one there.  Hence
\[
 \frac{y_n}{q^{a_n}c'_T}=q^m g(D)\in1+qA_0
                                      \subseteq1+q^{1/2}A_0.
 \tag{N-41}
\]
This proves the required residue condition.  The quotient by
the bar-lattice generator is
\[
 \frac{y_n}{q^{-a_n}c'_T}=q^{-m-2n}g(D),
 \tag{N-42}
\]
whose degree at infinity is $-2n\le0$.  Thus it is regular there.
Every $y_n$ lies in the actual Laurent form by (N-39), so all the
candidates belong to exactly the fibres specified in VIII.

Changing the initial vector on this line does not cure the defect.
For any $c'_T=h(q)d_T$ with $h\in K^\times$, put
$u=t/c'_T$ and $a_0=\operatorname{ord}_0(u)\in\tfrac12\mathbb Z$.
Multiplying $t$ by a nonzero rational constant normalizes the leading
coefficient of $q^{-a_0}u$ to one; this preserves (N-39) and bar
invariance.  For $a_n=a_0-n$ and $y_n=q^{-n}t$,
\[
 \frac{y_n}{q^{a_n}c'_T}=q^{-a_0}u\in1+q^{1/2}A_0,
 \qquad
 \frac{\bar y_n}{q^{a_n}c'_T}=q^{2n-a_0}u\in A_0.
 \tag{N-43}
\]
The second assertion is exactly $y_n\in\bar R^{a_n}$.
Thus the same family is admissible with every initial rational
normalization of the GT line.  Rational powers occurring here are
half-integers, so their powers of $q$ lie in $K=\mathbb Q(v)$.

To finish the argument, express the fixed nonzero $t$ in the ambient
KL tensor basis.  Every nonzero coefficient of $y_n$ has degree
equal to the corresponding degree of $t$ minus $n$, and the support
does not change.  More generally, every nonzero candidate in the
one-dimensional space (N-38) has this same support and differs from
$t$ by one nonzero coefficient-field element of finite valuation.  Its entire finite degree
tuple therefore differs by one finite scalar shift.  Taking $n$
sufficiently large makes the degree tuple of $y_n$ strictly smaller
at every nonzero coordinate.  It is then smaller under VIII's
partial order.  The reverse inequality fails: choose a minimal
nonzero coordinate in the ambient KL order.  Its old degree is
strictly larger, and no preceding coordinate has a smaller old
degree to satisfy the alternative clause in the definition.
The two tuples therefore compare strictly.  No candidate is a
global minimum.  The choice of the partial label order is irrelevant:
every finite nonempty poset has a minimal label, and its line has
the property just proved.

\paragraph{What this decides, and what a repair would change.}
The literal external minimum-degree prescription of VIII does not
construct $E(4)$, even after all GT matrices, path labels and initial
normalizations have been supplied.  Consequently there is no full
table singled out by that prescription whose signed constants,
coefficient field or dimension ring can be computed as requested.
The displayed rank-three examples and the one distinguished
rank-four element remain explicit elements and can be studied on
their own; this argument does not invalidate their algebraic formulas.

One possible repair is to require $y=\bar y$ explicitly.  That
excludes $q^{-n}t$ for $n>0$ because $\bar t=t\ne0$.  Another is
to bound or fix the permitted scaling exponents.  Neither repair
is part of (N-36), and neither is proved here to yield unique
minimum lifts or the printed examples.  VIII's section 4 gives a
different internal degree notion and explicitly begins its
discussion by assuming a bar-invariant element; our counterexample
does not refute that additionally restricted alternative.  The
text conjectures agreement of the two procedures, which cannot be
taken as established in the presence of the displayed external
defect.

Thus item 9(a) has a precise obstruction in its stated source
prescription, rather than a missing numerical computation.  A
repaired canonical-basis problem remains distinct from both the
proved integral nonfreeness in N-009 and the proved rational GT
construction in N-010.

\paragraph{Checks and literature boundary.}
The source passages checked are VIII's coefficient rings in
subsection 2.2, the whole of subsections 3.1--3.3, equation (13),
Conjectures 3.1 and 3.2, and the separate internal definition in
section 4.  The proof above is an elementary algebraic
counterexample, independent of a numerical experiment.  Its
rank-four extension uses N-010's independently reviewed rational
GT construction; the rank-one counterexample has no such input.
Targeted searches for the exact title with ``minimum'' and
``counterexample'', for ``Mulmuley minimum degree scaling canonical
basis'', and for ``0709.0751 Conjecture 3.2'' found no earlier report
of this defect.  The appropriate historical claim is limited to that
search and the supplied sources.  The plethysm branch independently
read the whole proof and the complete primary source passages;
its review is recorded in
\nolinkurl{nonstandard_plethysm_review_N012.md}.


\subsection{N-013: descent in an explicitly bar-fixed globalization variant}
\label{sec:nonstandard-bar-fixed-descent}

\paragraph{The variant and the claim.}
This subsection changes the domain of GCT VIII, subsection 3.3,
Conjecture 3.2, by explicitly imposing $y=\bar y$.  It does not
reinterpret that condition as part of the printed conjecture, which
is refuted under its displayed domain by N-012.  We work over
$K=\mathbb Q(v)$, put $q=v^2$, $\delta=v+v^{-1}$ and $D=\delta^2$,
and allow the scaling exponents $a\in\tfrac12\mathbb Z$, precisely
the rational powers $q^a$ that belong to $K$.

Let
\[
 R=\mathbb Q[v,v^{-1}],\qquad S=\mathbb Q[D],\qquad
 F=\mathbb Q(D),\qquad
 B_S=S\langle P_1,P_2,P_3\rangle,\qquad
 B_R=R\langle P_1,P_2,P_3\rangle.
 \tag{N-44}
\]
These are the actual subalgebras generated by the indicated words
inside the Hecke tensor algebra.  No saturation or localization
of $B_S$ or $B_R$ is taken.  Let $\{c'_T\}$ be the rational GT
matrix-unit basis supplied by N-010, and fix a genuine finite
partial order on its labels.  In particular, every $c'_T$ belongs
to $B_F=F\langle P_1,P_2,P_3\rangle$.  This is a specified rational
normalization and an additional choice of label order; it does
not claim that VIII uniquely prescribes either choice.

Let $A_0\subset K$ be the ring regular at $v=0$.  At an inductive
stage $T$, suppose that all lower scaling exponents $a_U$ have
been fixed.  Define
\[
 L_a=A_0q^ac'_T\oplus
       \bigoplus_{U<T}A_0q^{a_U}c'_U,
 \qquad
 \mathcal F_a=
 \{y\in B_R\cap L_a\cap\bar L_a:
       \bar y=y,\ y-q^ac'_T\in vL_a\}.
 \tag{N-45}
\]
The degree order is exactly VIII's external order on ambient KL
coefficient pole orders at $q=\infty$; the ambient coordinate
order can be any of the finite orders described there.

The following conclusions hold.
\begin{enumerate}
\item $B_S$ is a finite free $S$-module of rank $114$, and
      $B_R=R\otimes_S B_S$ inside the actual ambient algebra.
\item A global minimum over the union of the fibres
      $\mathcal F_a$ cannot have $a\in\tfrac12+\mathbb Z$.
\item If a global minimum exists, then at the same integral
      exponent there is a minimum in $B_S$.  In particular, if
      the minimizing pair $(a,e)$ is unique, then $a\in\mathbb Z$
      and $e\in B_S$.
\end{enumerate}
These statements do not prove existence or uniqueness of a
minimum.  Their consequence is conditional descent for this
specified variant, not a construction of VIII's printed $E(4)$.

\paragraph{The actual polynomial order is free.}
The ring $R$ is finite free of rank four over $S$.  Indeed
\[
 v^4+(2-D)v^2+1=0,\qquad
 v^{-1}=(D-2)v-v^3,
 \tag{N-46}
\]
so $1,v,v^2,v^3$ span.  They are independent over $F$: the four
automorphisms $v\mapsto v,v^{-1},-v,-v^{-1}$ of $K/F$ are distinct,
whereas (N-46) bounds $[K:F]$ by four.  Thus $[K:F]=4$, proving
the asserted independence over both $F$ and $S$.

The ambient integral Hecke tensor algebra, after replacing
integer coefficients by rational coefficients, is finite free
over $R$ and hence finite over $S$.  Since $S$ is noetherian,
its submodule $B_S$ is finitely generated.  It is torsion-free,
being a submodule of an $S$-torsion-free module.  The ring $S$
is a principal ideal domain, so $B_S$ is free.  Its rank is
$\dim_F B_F=114$ by the inherited dimension and $F$-form inputs
used in N-010.

Choose an $S$-basis $b_1,\ldots,b_{114}$ of $B_S$.  Every $P$-word
is an $S$-linear combination of these basis elements, so
$R\otimes_S B_S\to B_R$ is surjective.  The $b_j$ are also
$K$-independent, because $B_F$ is an $F$-form of the actual
generic algebra.  Therefore the map is injective, establishing
the first conclusion with no change to the word order.
This freeness over $\mathbb Q[D]$ is consistent with N-009:
that result concerns the different coefficient ring
$\mathbb Z[v,v^{-1}]$ at a mixed-characteristic maximal ideal.

\paragraph{Two commuting semilinear involutions.}
Let bar be the usual Hecke involution, with $\bar v=v^{-1}$
and $\overline{T_i}=T_i^{-1}$.  Every $P_i$ and every element
of $S$ is fixed by bar.  Define a second semilinear involution
$\beta$ by
\[
             \beta(v)=-v,\qquad \beta(T_i)=-T_i.
 \tag{N-47}
\]
The quadratic relation is preserved because both $T_i$ and
$v-v^{-1}$ change sign; each braid relation is homogeneous
and has the same sign on its two sides.  Hence (N-47) defines
an involution of the Hecke algebra and its tensor product.
It commutes with bar.  Since $C'_i=T_i+v^{-1}$ and
$C_i=T_i-v$ both change sign, $P_i=C'_i\otimes C'_i+C_i\otimes C_i$
is fixed by $\beta$.  Also $\beta(D)=D$.

Consequently every $b_j$ and every $c'_T\in B_F$ is fixed
by both involutions.  In particular, $\beta$ does not permute
the GT lines or their labels.  It preserves every span cut out
by the chosen label order.  This follows from the actual
$F$-word expressions, not from an assumed symmetry of the
abstract branching diagram.

The invariant Laurent ring is $R^{\mathrm{bar}}=\mathbb Q[\delta]$:
each symmetric Laurent polynomial is a polynomial in
$v+v^{-1}$, by the recurrence
$v^{n+1}+v^{-n-1}=\delta(v^n+v^{-n})-(v^{n-1}+v^{-n+1})$.
Using the independent $b_j$ and $\beta(\delta)=-\delta$ gives
\[
 B_R^{\mathrm{bar}}=B_S\oplus\delta B_S,\qquad
 (B_R^{\mathrm{bar}})^{\beta=+1}=B_S,\qquad
 (B_R^{\mathrm{bar}})^{\beta=-1}=\delta B_S.
 \tag{N-48}
\]
These are identities in the actual order, so division by
$\delta$ on the last summand remains inside $B_R$.

\paragraph{The effect on degrees and on the fibres.}
In either usual KL convention,
$\beta(C_w)=(-1)^{\ell(w)}C_w$.  For example, this follows
from the standard expression for $C_w$ in which the exponent
of every coefficient of $T_x$ is congruent to
$\ell(w)-\ell(x)$ modulo two; replacing $v$ and every $T_i$
by their negatives gives the displayed sign.  Thus every
ambient KL tensor coordinate of $\beta y$ is the corresponding
coordinate of $y$, evaluated at $-v$ and multiplied by a sign.
It follows that
\[
 d(\beta y,z)=d(y,z),\qquad
 d(y_1+y_2,z)\le\max\{d(y_1,z),d(y_2,z)\}.
 \tag{N-49}
\]
The convention $d(0,z)=-\infty$ is used here and below.

Fix $a\in\tfrac12\mathbb Z$ and set $s=(-1)^{2a}$.
Each generator $q^{a_U}c'_U$ of the lower lattice and
$q^ac'_T$ is an eigenvector of $\beta$ with eigenvalue $+1$
or $-1$.  The coefficient ring $A_0$ is preserved by $\beta$,
and its residue field is $\mathbb Q$, on which $\beta$ is the
identity.  Therefore, if $y\in\mathcal F_a$, then
$s\beta(y)\in\mathcal F_a$.  The vector
\[
                   y_s=\tfrac12(y+s\beta(y))
 \tag{N-50}
\]
also belongs to $\mathcal F_a$, is nonzero by its specified
residue, and satisfies $\beta(y_s)=s y_s$.
Its degrees are coordinatewise at most those of $y$ by
(N-49).  Averaging preserves membership in the actual order
because its coefficient ring is rational; no division by two
in the integral order of N-009 is asserted.

If $s=-1$, (N-48) permits $z=y_s/\delta\in B_S$.
In the $T$ coordinate, division by $\delta$ changes the
admissible exponent from $a$ to $a+1/2$, since
\[
 \frac{q^a/\delta}{q^{a+1/2}}=\frac1{1+q}\in1+vA_0.
 \tag{N-51}
\]
In every lower coordinate, the factor
$1/\delta=v/(1+v^2)$ increases its valuation at zero by one
in the $v$ variable.  A coefficient already in $vA_0$
therefore remains in $vA_0$.  Hence $z\in L_{a+1/2}$ has the
required residue.  It is bar fixed, so membership in
$\bar L_{a+1/2}$ follows by applying bar to its membership
in $L_{a+1/2}$.  Thus
\[
                 y_s/\delta\in\mathcal F_{a+1/2}.
 \tag{N-52}
\]
At every coordinate where $y_s$ is nonzero, division by
$\delta$ decreases its degree at $q=\infty$ by $1/2$.
At every coordinate where $y$ is nonzero but $y_s$ is zero,
the new degree is $-\infty$.  Therefore (N-49)--(N-52)
strictly decrease the degree at every nonzero coordinate of
$y$, and never introduce a new nonzero coordinate.

This gives strict decrease in VIII's partial degree order,
not merely in a total refinement.  At a minimal coordinate
of the nonzero support of $y$, the reverse comparison fails:
the new degree is strictly smaller and all preceding
coordinates have value $-\infty$ for both vectors.
Thus no candidate with $s=-1$ can be a global minimum.

If $s=+1$ and $y$ is a global minimum, (N-50) gives a vector
$y_+\in B_S\cap\mathcal F_a$ with coordinatewise no larger
degree.  It is also a global minimum: for any competing
candidate $w$, each comparison witnessing $d(y)\le d(w)$
still holds with $y_+$ in place of $y$.  Indeed a coordinate
inequality remains true after lowering that coordinate, and
a strictly smaller preceding coordinate remains strictly
smaller after such lowering.  If the global minimum is
unique, $y=y_+\in B_S$.  This proves all three assertions.

\paragraph{What the descent does and does not settle.}
For the stated rational normalization, a unique bar-fixed
external globalization, if it exists, does not require a
coefficient extension of $\mathbb Q(D)$ at rank four.
This conclusion is stronger than merely splitting the generic
algebra: the resulting vectors would lie in its actual
$\mathbb Q[D]$ word order.  Their multiplication coefficients
in that rational basis would then belong to $F$, although
finite Laurent-polynomial or signed-positivity properties
still require separate proofs, as N-007 demonstrates.

There is also a concrete finite-linear-algebra consequence.
Use the flip-symmetric ambient KL coordinates and write a
$\beta$-fixed, bar-fixed coordinate as
$\delta^{\epsilon_z}p_z(D)$, where
$\epsilon_z\equiv\ell(w_1)+\ell(w_2)\pmod2$ and
$p_z\in\mathbb Q[D]$.  A fixed finite upper bound $h_z$
on every nonzero ambient degree permits only
$\deg p_z\le h_z-\epsilon_z/2$.  Within that finite-dimensional
space, actual-order membership, the lower-label support
conditions, and any fixed-scaling residue conditions are
linear conditions over $\mathbb Q$ once a polynomial module
basis is supplied.  This yields an exact feasibility problem
for every bounded box.

It does not, by itself, replace the unbounded optimization by
one bounded box.  A strict improvement at an earlier ambient
coordinate can permit arbitrarily large later coordinates in
VIII's partial degree order.  Nor does it select an order for
modules whose maximal $H$-weights are incomparable.  These
existence, uniqueness, order and global-cutoff issues remain
explicit gates for the repaired construction.


\subsection{N-014: the printed highest-weight label order is not defined}
\label{sec:nonstandard-label-order}

\paragraph{The precise source rule.}
GCT VIII, subsection 2.3, printed p.~11
(arXiv:0709.0751v2), defines $\mu(\alpha)$ to be the highest
$H_q$-weight in the restriction of $W_{q,\alpha}$ and sets
$\alpha\le\alpha'$ if and only if
$\mu(\alpha)\le\mu(\alpha')$ in the usual root-cone order.
The text permits several highest-weight vectors, but does not
specify a set of incomparable highest weights or a tie rule
for distinct labels with the same highest weight.
Subsection 3.3, step (1), printed p.~20, transfers this order
to the nonstandard Hecke labels by duality; step (2) then
orders GT tableaux lexicographically.

There are two distinct problems with this literal definition
already among the rank-four modules in the supplied examples.
First, the required single greatest weight need not exist.
Second, even where it exists, distinct labels can have the same
weight, so the stated ``if and only if'' gives a preorder,
not a partial order on those labels.  The argument below
uses exact modules, rather than relying on the floating-point
qualification in VII's experimental table.

\paragraph{An actual module with two incomparable maximal weights.}
Let $M_\lambda$ denote the ordinary Hecke simple of shape
$\lambda\vdash4$.  Write $N=M_4\otimes M_{22}$, one of the
two-dimensional simples proved distinct from all other factors
in N-010.  For $H_q=GL_q(2)\times GL_q(2)$, let $V_\lambda$
be the standard quantum Weyl module of shape $\lambda$.
Ordinary quantum Schur--Weyl duality and N-010 give the exact
$H_q$-module identity
\[
 \operatorname{Hom}_{B_{4,K}}(N,X_q^{\otimes4})
       \cong V_4\otimes V_{22}\ \oplus\ V_{22}\otimes V_4.
 \tag{N-53}
\]
Here the Hom denotes the multiplicity space for the actual
nonstandard action on the tensor power, and $K=\mathbb Q(v)$.
In the duality terminology of VII this is $W_2$.

For completeness, the needed exact tensor restrictions for the
three ordinary shapes with at most two rows are
\[
\begin{array}{c|c}
\text{ordinary tensor product}&\text{nonstandard constituents}\\\hline
M_4\otimes M_4&\epsilon_+\\
M_4\otimes M_{31}&M_4\otimes M_{31}\\
M_4\otimes M_{22}&N\\
M_{31}\otimes M_{31}&\epsilon_+\oplus\Lambda^2M_{31}\oplus S'M_{31}\\
M_{31}\otimes M_{22}&M_{31}\otimes M_{22}\\
M_{22}\otimes M_{22}&\epsilon_+\oplus\epsilon_-\oplus S'M_{22}.
\end{array}
 \tag{N-54}
\]
Tensor flip gives the reversed pairs.  The symmetric and
alternating decompositions, the simultaneous $D$-eigenline,
the quotient models and pairwise distinctness of all listed
simples were verified exactly in N-010.  The alternating
square of $M_{22}$ has $P_i=0$: its dimension is one and a
rank-one matrix $C'_i$ has tensor eigenvalues yielding zero
on that alternating line.  Thus $N$ occurs once in each of
$M_4\otimes M_{22}$ and its flip, and nowhere else in (N-54).
This proves (N-53) without a numerical specialization and
without assuming the entire nonstandard duality conjecture.
If that conjecture is used to name the paired quantum-group
module, its restriction must be this multiplicity space.

The two highest weights in (N-53) are
\[
 \lambda=((4,0),(2,2)),\qquad
 \mu=((2,2),(4,0)).
 \tag{N-55}
\]
Their difference is $(2,-2;-2,2)$.  The simple-root cone of
$GL_2\times GL_2$ consists of
$(a,-a;b,-b)$ with $a,b\ge0$.  Neither difference belongs
to this cone.  Each weight in (N-55) is maximal in the
direct sum, and they are incomparable.  Therefore there is
no single highest weight $\mu(N)$ in the sense required
by VIII's displayed order rule.

This is exactly the decomposition displayed in VII,
subsection 7.2.2, printed p.~48 (arXiv:0709.0749v2).
The accompanying computational qualification is on printed
p.~49.  Equations (N-53)--(N-54) supply an exact independent
justification for the particular decomposition used here.

\paragraph{The defect persists for a faithful full-rank realization.}
The preceding example is not repaired by increasing the
dimensions of the two standard vector spaces so that all
ordinary Hecke shapes occur.  To see this, the involution
$\theta(T_i)=-T_i^{-1}$ conjugates ordinary Specht shapes
and twisting one tensor factor sends $P_i$ to $D-P_i$.
Thus the remaining shapes $211$ and $1111$ are obtained
from $31$ and $4$ by the same twist.  The shape $22$ is
self-conjugate.  On the explicit $N$ matrices of N-010,
the intertwiner
\[
       J=\begin{pmatrix}1&1\\-D&-1\end{pmatrix},\qquad
       \det J=D-1,
       \qquad J P_i=(D I-P_i)J
 \tag{N-56}
\]
verifies directly that $N$ is invariant up to isomorphism
under this nonstandard twist.

Twisting both ordinary tensor factors leaves the nonstandard
action unchanged.  Twisting one factor preserves dimensions
and simplicity of constituents.  Consequently the products
involving a shape $31$ or $211$ have constituents of
dimensions $1,3,5$ or the single dimension $6$, as applicable,
and cannot add $N$ to their decompositions.  The product
$M_{22}\otimes M_{22}$ is itself invariant under the twist;
its one-dimensional constituents are interchanged, so its
unique two-dimensional constituent $S'M_{22}$ is invariant
and remains distinct from $N$.
It follows that, among all five shapes of size four, $N$
occurs exactly in the four ordered pairs
\[
          (4,22),\quad(22,4),\quad(1111,22),\quad(22,1111).
 \tag{N-57}
\]
Each occurs with multiplicity one.  For $GL_q(d)\times GL_q(e)$,
the last two summands are omitted when the relevant
four-row Weyl module vanishes.  When present their highest
weights are dominated by those of the corresponding first
two summands.  The maximal antichain remains
$\{(4,22),(22,4)\}$ for every $d,e\ge2$.
Thus the failure of a greatest weight applies also when
the full rank-$114$ algebra acts faithfully.

\paragraph{Equal highest weights are a separate issue.}
For $GL_q(2)\times GL_q(2)$, (N-54) also proves
\[
 W_{\epsilon_-}\cong W_{S'M_{22}}
          \cong V_{22}\otimes V_{22},\qquad
 W_{\Lambda^2M_{31}}\cong W_{S'M_{31}}
          \cong V_{31}\otimes V_{31}
 \tag{N-58}
\]
as $H_q$-modules.  The respective nonstandard simple labels
are distinct by N-010.  VII's notation for these pairs is
$(W_1,W_6)$ and $(W_3,W_7)$, respectively, in subsection 7.2.2.
The literal comparison through their equal highest weights
therefore relates each pair in both directions, violating
antisymmetry on labels.  One may quotient the label set by
this equivalence, but then the inductive subspace attached
to a minimal type need not be a single GT line.  That is a
different prescription, not a consequence of the stated
partial-order definition.

\paragraph{Explicit alternative orders, and the remaining boundary.}
One possible repair assigns to every label its finite
antichain $A_\alpha$ of maximal $H_q$-weights and compares
these antichains by
\[
 A\preceq B\quad\Longleftrightarrow\quad
 \text{for each }a\in A\text{ there exists }b\in B
                       \text{ with }a\le b.
 \tag{N-59}
\]
This is a partial order on antichains: equivalently, it is
inclusion of their generated lower sets.  A partial order on
the original labels is obtained by using strict inclusion
for distinct antichains and declaring distinct labels with
equal antichains incomparable.  A total extension may instead
resolve every tie by an explicitly recorded additional rule.
For example, either $\epsilon_-<S'M_{22}$ or its reverse
is possible after replacing equality of their antichains
by a chosen tie order; likewise either order is possible
for $\Lambda^2M_{31}$ and $S'M_{31}$ in the two-row example.
These choices alter the lower-label spans in the inductive
globalization and therefore cannot be left implicit.

Another possible convention chooses one maximal weight
according to a fixed total refinement of the product root
order, and then separately resolves equal weights.  For
(N-55), giving the first factor priority chooses $\lambda$;
giving the second factor priority chooses $\mu$.  Those are
two different extra choices, neither specified by VIII.

No theorem here says that a repaired order cannot produce a
useful basis.  The conclusion is the precise source-level
one: the rank-four GT data in N-010 do not complete VIII's
printed label-order recipe.  Together with N-012, this
explains why computing a table cannot by itself identify
that table with a uniquely prescribed $E(4)$.


\subsection{N-015: a certified basis of the actual $\mathbb Q[D]$ order}
\label{sec:nonstandard-polynomial-order}

\paragraph{Exact result.}
The actual polynomial word order
$B_S=\mathbb Q[D]\langle P_1,P_2,P_3\rangle$ has a basis of
114 elements over $S=\mathbb Q[D]$.  The accompanying certificate
specifies every element in 300 ambient coordinates and supplies
a straight-line expression for it in 148 actual $P$-words of
length at most ten.  Exact multiplication by all three
generators preserves the resulting module.  Thus the certificate
describes the full unsaturated order, in every word degree.

The 114 basis rows have distinct pivot coordinates.  Their
monic pivot degrees have distribution
\[
\begin{array}{c|rrrrrrr}
\text{degree}&0&1&2&3&4&5&6\\\hline
\text{number of pivots}&4&36&39&18&12&2&3.
\end{array}
 \tag{N-60}
\]
Every coordinate polynomial of every final row has degree at
most eight.  This gives an explicit integral basis over
$\mathbb Q[D]$, not over $\mathbb Z[v,v^{-1}]$; the latter
is excluded by N-009.  It also independently proves the
generic dimension 114 and the $\mathbb Q(D)$-form property
used as attributed inputs earlier in N-010.

\paragraph{Ambient coordinates.}
Let $C_w$ be the opposite KL basis of the ordinary rank-four
Hecke algebra, with $C_s=T_s-v$.  Order permutations
lexicographically in one-line notation, and index coordinates
by the 300 unordered pairs $a\le b$ of the 24 permutations.
Put $\epsilon_{a,b}=\ell(a)+\ell(b)\pmod2$.  The coordinate
vector $f_{a,b}(D)$ denotes the coefficient of
\[
 \delta^{\epsilon_{a,b}}\begin{cases}
 C_a\otimes C_a,&a=b,\\
 C_a\otimes C_b+C_b\otimes C_a,&a<b.
 \end{cases}
 \tag{N-61}
\]
In particular an off-diagonal coordinate is the coefficient
of either ordered tensor term, not twice that coefficient.
The tensors in (N-61) are independent over $K=\mathbb Q(v)$.
Every $P$-word is fixed by tensor flip, bar, and $\beta$ of
N-013, so its coordinates are polynomials in $D$.  This last
fact also follows directly from the exact generator matrices
used below, with polynomial initial vector $1$.

The first construction computes the right action of $C_s$
on each opposite KL vector, then uses
\[
 P_s=2C_s\otimes C_s+
       \delta(C_s\otimes1+1\otimes C_s)+D\,1.
 \tag{N-62}
\]
The parity factors in (N-61) make every resulting entry an
integer polynomial in $D$.  The script asserts that each
power of $\delta$ left after this change of coordinates is
nonnegative and even.  It constructs all 900 generator
columns, rather than evaluating at a parameter.

\paragraph{The polynomial module computation.}
Starting with $1$, the construction tests products of accepted
words by the three generators.  It inserts any vector outside
the current polynomial row module using Euclidean elimination
over the PID $S$.  At a shared pivot with coefficients $p,a$,
write $g=sp+ta=\gcd(p,a)$ and replace the two rows $b,x$ by
\[
           sb+tx,\qquad (p/g)x-(a/g)b.
 \tag{N-63}
\]
The transformation has determinant
$s(p/g)+t(a/g)=1$, so it preserves the actual $S$-span.
Normalizing a pivot divides only by a nonzero rational
constant.  Subtracting polynomial multiples of a pivot row
reduces entries above it.  No polynomial divisor is inverted.

The process accepts 148 words and stops after their complete
generator closure is contained in the row module.  There
are 316 pivot insertions or gcd updates.  A search stopping
at word length ten alone would not prove this assertion;
the proof uses the exact closure verification after the
last accepted word.

Every row operation is stored as a node of a straight-line
program, whose leaves are the accepted words and whose
internal nodes are sums of two preceding nodes multiplied
by specified polynomials in $S$.  Removing unused nodes
leaves 34,026 nodes.  This avoids the very large expanded
word coefficients produced by an initial, less compact
certificate format.  All final rows were reconstructed
exactly from these nodes.  Thus every final row is in the
actual word order.

\paragraph{Independent arithmetic replay and all-degree proof.}
The separate script
\nolinkurl{verify_polynomial_order_rank4.py} imports neither
module-search script.  It reconstructs $C_w$ in the standard
$T$ basis by the semilinear substitution
$v\mapsto v^{-1}$, $T_i\mapsto-T_i$, which sends
$C'_w$ to $(-1)^{\ell(w)}C_w$.  Triangular standard-basis
elimination then gives the complete right action of each
$T_s$.  It constructs the generator columns using instead
\[
 P_s=2T_s\otimes T_s-(v-v^{-1})(T_s\otimes1+1\otimes T_s)
                         +(v^2+v^{-2})\,1.
 \tag{N-64}
\]
Every Laurent polynomial is converted to the form in (N-61)
by repeatedly subtracting its leading multiple of a power
of $\delta$, checking all coefficients and its bar symmetry.
This verifies the full polynomial operators through a
different Hecke arithmetic path.

From those operators the verifier reconstructs all 148
accepted words and replays all 34,026 straight-line nodes.
It verifies all final rows exactly, their distinct monic
pivots, membership of $1$, membership of all listed words,
and all 342 products of a final basis row by a generator.
Membership uses ordinary triangular polynomial division:
a nonzero remainder at a pivot rejects membership.  This
does not test membership in a saturation or in the generic
fraction-field span.

Let $M$ be the $S$-span of the certified rows.  The
straight-line expressions prove $M\subseteq B_S$.
The conditions $1\in M$ and $MP_i\subseteq M$ imply by
induction that every $P$-word lies in $M$, so $B_S\subseteq M$.
Consequently
\[
                         M=B_S.
 \tag{N-65}
\]
Distinct nonzero pivot coordinates imply linear independence
over $S$, over $F=\mathbb Q(D)$, and over $K$.  Therefore
the rows are an actual polynomial basis of $B_S$, and the
generic $K$-span of every word has dimension exactly 114.
Their $F$-span extends to that same generic algebra over $K$.
This proves the dimension and form assertions without
assuming either one as an upper bound.

\paragraph{Certificates, scope, and checks.}
The files
\nolinkurl{polynomial_order_rank4_verified.json} and
\nolinkurl{polynomial_order_rank4_rows.json.gz} contain the
word list, all coordinates, the straight-line program and
the exact summary.  The compressed certificate is 972,422
bytes.  The search script
\nolinkurl{polynomial_order_rank4_flint.py} requires
\texttt{python-flint}; the first slower Euclidean prototype
\nolinkurl{polynomial_order_rank4.py} is retained as search
history.  The compact exact construction completed in
about 28 seconds in this run; the independent arithmetic
replay completed in about 14 seconds.  These elapsed times
are observations, not complexity bounds.

The independent arithmetic replay was authored within the
same branch.  It is not labelled an independent whole-proof
review by another agent.  The separate event ledger records
such reviews when received.  No minimum-degree globalization,
signed structure-constant table, or canonical normalization
is inferred from the polynomial order computation itself.


\subsection{N-016: a finite global decision procedure for the bar-fixed variant}
\label{sec:nonstandard-finite-cutoff}

\paragraph{Statement and boundary.}
Keep the explicitly modified domain, coefficient field,
rational GT normalization and chosen genuine finite label
order of N-013.  At a fixed inductive stage $T$, assume that
all lower scaling exponents $a_U\in\tfrac12\mathbb Z$ have
already been fixed.  There is an effective finite procedure
which decides whether a global minimum over all current
scaling exponents exists, whether the minimizing pair is
unique, and, if unique, computes it exactly.

The procedure uses the actual $S=\mathbb Q[D]$ order, not
its saturation in an ambient polynomial module.  Its inputs
are a polynomial basis of that order, its rational GT
coordinate map, the label and ambient orders, and the lower
exponents.  N-010 and N-015 supply the algebraic data at rank
four.  Polynomial Smith or Hermite operations and finite
linear algebra over $\mathbb Q$ suffice.  No efficient bound
on the running time is asserted, and this theorem does not
assert that every stage has a unique minimum.  In particular,
it is not a computed full basis or a table of structure
constants.

\paragraph{An exact finite-dimensional quotient.}
Write $c'_U$ for the fixed GT vectors, $B_S$ for the actual
polynomial word order, and
\[
 E_T=\bigoplus_{U\le T}F c'_U,\qquad
 M_T=B_S\cap E_T,
 \qquad k=\#\{U:U\le T\}.
 \tag{N-66}
\]
If $h_U(y)\in F$ is the $U$-th GT coordinate, put
\[
 b_U=\left\lfloor-a_U-\tfrac12\right\rfloor\quad(U<T),
 \qquad
 V=\{y\in M_T:\deg_D h_U(y)\le b_U\text{ for every }U<T\}.
 \tag{N-67}
\]
Here $\deg_D$ of a rational function is numerator degree
minus denominator degree, and $\deg_D0=-\infty$.
Since $D=q^{-1}(1+q)^2$, its negative is the valuation at
$q=0$.  Hence the inequalities in (N-67) are exactly the
lower-coordinate zero-residue conditions in (N-45), for
an element of $B_S$.  The factor $1/2$ is essential because
the residue quotient is by $v=q^{1/2}$.

The module $M_T$ is finite free over $S$ and has rank $k$.
Indeed it is an $S$-submodule of the finite free module
$B_S$, and clearing the finitely many denominators of any
$F$-linear expression for $c'_U$ shows that its $F$-span is
exactly $E_T$.  Let
\[
 I_T=B_S\cap F c'_T=S t,
 \tag{N-68}
\]
where $t\ne0$ is a generator of this actual rank-one
submodule.  A generator exists because $S$ is a PID.
The lower-coordinate projection
$\pi:M_T\to F^{k-1}$ has kernel $S t$.

Choose an $S$-basis $u_1,\ldots,u_k$ of $M_T$ and a nonzero
common denominator $\Delta(D)$ for all lower GT coordinates
of these vectors.  Then
\[
 N_0=\Delta\pi(M_T)\subset S^{k-1}
 \tag{N-69}
\]
is an actual finitely generated polynomial submodule.
For $y\in V$, the numerator vector
$n=\Delta\pi(y)$ lies in
\[
 W_0=\bigoplus_{U<T}\mathbb Q[D]_{\le b_U+\deg\Delta},
 \tag{N-70}
\]
where a polynomial space with negative upper bound is zero.
The kernel of $V\to N_0\cap W_0$ is $S t$, and the map is
surjective by the definitions of $N_0$ and $V$.
Consequently
\[
 V/(S t)\cong N_0\cap W_0,
 \qquad
 \dim_{\mathbb Q}V/(S t)
 \le\sum_{U<T}\max\{0,b_U+\deg\Delta+1\}.
 \tag{N-71}
\]
This proves finite-dimensionality with a concrete bound.
It does not replace $N_0$ by $S^{k-1}$.

Here is an explicit way to compute the quotient and its lifts.
First find $M_T$ as the polynomial kernel of the excluded
GT-coordinate rows in an $S$-basis of $B_S$, after clearing
their denominators.  This computes the actual kernel inside
$B_S$.  Let $A$ be the polynomial matrix with columns
$\Delta\pi(u_j)$.  It has rank $k-1$ over $F$.
A polynomial Smith reduction
$U A W=\operatorname{diag}(d_1,\ldots,d_{k-1},0)$, with
$U,W$ unimodular over $S$, determines membership in $N_0$:
the $i$-th coordinate of $Un$ must be divisible by $d_i$.
For $n\in W_0$, the remainders on division by these fixed
polynomials are finite linear conditions on its rational
coefficients.  Solve those conditions to obtain a
$\mathbb Q$-basis $n_1,\ldots,n_m$ of $N_0\cap W_0$.
The Smith operations lift each $n_i$ to a polynomial vector
of coefficients in the $u_j$, giving $y_i\in V$.
The remaining column of $W$ gives a generator $t$ of the
kernel.  Every division used to lift is performed only
after the divisibility condition has been checked.
Thus this construction uses neither saturation nor an
unrecorded localization.

We have now obtained the exact decomposition, as a
$\mathbb Q$-vector space,
\[
                    V=S t\oplus
                          \operatorname{span}_{\mathbb Q}\{y_1,\ldots,y_m\}.
 \tag{N-72}
\]
For $k=1$ the lower projection is empty, $m=0$, and the
same assertion reads simply $V=S t$.

\paragraph{An explicit global degree cutoff.}
Use the ambient coordinates (N-61).  Write
\[
 t_z=\delta^{\epsilon_z}r_z(D),\qquad
 (y_i)_z=\delta^{\epsilon_z}p_{i,z}(D).
\]
Define the nonnegative integer
\[
 B=\max\left(\{0\}\cup
       \{\deg p_{i,z}-\deg r_z:
                   r_z\ne0,\ p_{i,z}\ne0\}\right).
 \tag{N-73}
\]
This is an explicitly computable maximum of finitely many
degrees of polynomials already constructed above.  Every
$y\in V$ has a unique expression
$y=p(D)t+\sum_i c_i y_i$, with $p\in S$ and $c_i\in\mathbb Q$.
If $\deg p>B$, then at every coordinate with $t_z\ne0$
the term $p(D)t_z$ has strictly higher degree than every
$(y_i)_z$, and also strictly higher degree than $t_z$.
There is therefore no leading cancellation, and
\[
             d(y,z)=\deg p+d(t,z)>d(t,z)
                       \qquad(t_z\ne0).
\]
At every coordinate with $t_z=0$, its degree is $-\infty$,
which is at most the corresponding degree of $y$.
In fact the normalized $t$ strictly lowers the degree at
every coordinate where $y$ is nonzero, and its support is
contained in that of $y$.

To check admissibility, put $g(D)=h_T(t)=t/c'_T\in F^\times$,
let $\kappa\in\mathbb Q^\times$ be its leading coefficient
at $D=\infty$, and replace $t$ by $t/\kappa$.
It remains a generator of the same actual module.
With
\[
                     a_0=-\deg_D g\in\mathbb Z,
\]
the ratio $h_T(t)/q^{a_0}$ is one modulo $vA_0$.
Every lower GT coordinate of $t$ is exactly zero, so all
previous lower exponents and residue conditions are retained.
It is bar fixed and lies in the actual order.  It therefore
belongs to the repaired fibre $\mathcal F_{a_0}$.

This comparison is over all current exponents, precisely
the variable-scaling problem.  It does not claim that $t$
belongs to the fibre of the original candidate's arbitrary
current exponent.  The previously fixed lower fibres are
unchanged.  Since the degree improvement is coordinatewise
and strict on the old nonzero support, it is strict in
VIII's partial degree order as well: a minimal coordinate
of that support prevents the reverse comparison.
Thus no global minimum can have $\deg p>B$ in (N-72).

All polynomial candidates that can minimize consequently
belong to the explicit finite-dimensional space
\[
 V_B=\operatorname{span}_{\mathbb Q}
       \{y_1,\ldots,y_m,t,Dt,\ldots,D^B t\}.
 \tag{N-74}
\]
At each ambient coordinate a concrete upper degree bound is
\[
 H_z=\max\left(\{\deg p_{i,z}:p_{i,z}\ne0\}
           \cup\{B+\deg r_z:r_z\ne0\}\right),
 \qquad
 d(y,z)\le H_z+\epsilon_z/2,
 \tag{N-75}
\]
with an empty maximum interpreted as $-\infty$.

\paragraph{Finite feasibility and the global partial order.}
The current-coordinate functional $h_T$ on $V_B$ is rational
in $D$ with a fixed common denominator.  Multiplying by that
denominator gives a polynomial numerator whose coefficients
are linear in the rational coordinates of $y\in V_B$.
There are only finitely many possible degrees of a nonzero
numerator.  For each one, requiring the leading coefficient
of $h_T$ to be one and all higher numerator coefficients to
vanish is an affine linear system.  Its current scaling
exponent is exactly $a=-\deg_D h_T$.  All lower constraints
hold already by the definition of $V_B$.

Similarly, fixing upper degree bounds on ambient coordinates
is a finite system of linear coefficient vanishings.
Equations (N-74)--(N-75) leave only finitely many possible
ambient degree profiles.  Their exact attainability is
decidable: in the relevant affine solution space, impose
the upper vanishings and require each designated leading
coefficient to be nonzero.  Over the infinite field
$\mathbb Q$, a finite collection of such nonvanishing
requirements is simultaneously attainable if and only if
none of the corresponding affine functionals vanishes
identically on the solution space.  This follows because
a finite union of proper affine hyperplanes cannot cover
that space.  A rational point avoiding them can be found
by a finite search once their properness is established.

One must compare these finitely many attainable profiles
using VIII's actual partial degree order.  Minimizing a
single lexicographic refinement is insufficient.  The
source comparison is itself a partial order: it is the
intersection of the lexicographic orders over all linear
extensions of the finite ambient-coordinate order.  In
one direction, a first differing coordinate with the wrong
sign would need an earlier strict improvement.  Conversely,
if the source condition fails at a coordinate $z$, its lower
set contains no strict improvement; a linear extension
listing that lower set first makes the lexicographic
comparison fail.  This proves the equivalence.

Test whether any attainable profile is at most every other
attainable profile under this comparison.  If none is,
there is no global minimum, even though minimal profiles
exist.  If one is, its realizing elements are global
minima among the polynomial candidates.  Every polynomial
candidate outside $V_B$ is strictly improved by the
normalized $t\in V_B$, so this finite test is global over
all degrees and all current exponents.

\paragraph{Returning to all bar-fixed candidates and testing uniqueness.}
The reduction to $B_S$ loses no global-minimum existence
information.  For every candidate in $B_R^{\mathrm{bar}}$,
N-013 supplies a candidate in $B_S$ with coordinatewise
no larger degree: parity averaging suffices at an integral
current exponent; at an odd half-integer exponent one also
divides the odd projection by $\delta$.  If the resulting
polynomial candidate is outside $V_B$, the cutoff replaces
it by the normalized $t$.  Therefore every candidate in
the full repaired domain is dominated by some candidate
from the finite space $V_B$.

It follows that the full bar-fixed domain has a global
minimum if and only if the polynomial finite test does.
Any full-domain minimum has the same degree profile as
its polynomial projection: a strict coordinatewise
decrease would contradict global minimality.  Its profile
is consequently bounded by (N-75).

To test uniqueness of the pair $(a,y)$ in the full domain,
work within that finite ambient box.  Every bar-fixed
Laurent coefficient is a polynomial in $\delta$; the bound
at coordinate $z$ permits degree in $\delta$ at most
$2H_z+\epsilon_z$, or requires zero when $H_z=-\infty$.
Membership in the actual bar-fixed order is computed using
the decomposition $B_R^{\mathrm{bar}}=B_S\oplus\delta B_S$
of N-013 and the polynomial-module membership tests.
Explicitly, test $y_+=(y+\beta y)/2\in B_S$ and
$y_-/\delta=(y-\beta y)/(2\delta)\in B_S$ separately.
In the normalized coordinates (N-61), the second test can
require a divisibility check by $D$ before its coordinates
are polynomials; that condition is part of actual module
membership and is not discarded.
The GT support restrictions and the prescribed residue
conditions are finite linear conditions after multiplying
by fixed common denominators.  Only finitely many possible
current leading valuations occur in this bounded space.
Odd half-integer current exponents cannot minimize by N-013.

For each remaining exponent, solve the affine system with
ambient degree bounded above by the already found least
profile.  Every solution realizes that profile: a strictly
smaller coordinatewise profile would contradict its least
property.  A positive-dimensional affine solution space
gives nonuniqueness, as do solutions for two distinct
current exponents.  Exactly one point at exactly one
exponent gives the unique minimizing pair.  This includes
possible odd-parity homogeneous perturbations, rather
than inferring full-domain uniqueness solely from the
polynomial subspace.

\paragraph{A direct linear implementation without profile enumeration.}
There is a smaller exact implementation of the same finite
decision.  On any subspace $W\subseteq V_B$, a candidate
exists if and only if $h_T$ does not vanish identically on
$W$.  Indeed every vector with $h_T(y)\ne0$ can be divided
by its nonzero rational leading coefficient and assigned
the exponent $a=-\deg_D h_T(y)$; the lower bounds are
homogeneous and are unaffected.  Testing nonvanishing of
this rational functional on $W$ is finite rational linear
algebra after clearing its fixed denominator.

Choose one linear extension of the ambient-coordinate order.
Starting with $V_B$, minimize its successive coordinate
degrees in that total order.  At a coordinate $z$, try zero
first and then degrees $0,1,\ldots,H_z$ for its polynomial
part.  Impose the corresponding linear coefficient
vanishings, retaining the first bound for which $h_T$ is
not identically zero on the remaining subspace.  The process
terminates and produces the lexicographically least degree
profile $d^*$ and a nonzero subspace realizing it.  Every
vector in that subspace with $h_T\ne0$, after normalization,
has exactly that profile; a smaller degree at any first
differing coordinate would contradict the greedy step.

This one lexicographic minimum is only a candidate for the
global partial-order minimum.  For each coordinate $z$ with
$d^*_z\ne-\infty$, test whether there exists $y\in V_B$,
$h_T(y)\ne0$, satisfying
\[
 d(y,z)<d^*_z,\qquad
 d(y,x)\le d^*_x\quad\text{for every strict predecessor }x<z.
 \tag{N-76}
\]
There is no further bound beyond the already fixed finite
space $V_B$.  Strictness in (N-76) lowers the integer
polynomial degree bound by one, so this is again a
homogeneous linear-system test together with nonvanishing
of $h_T$.

Condition (N-76) is precisely a witness that the defining
comparison $d^*\le d(y)$ fails at $z$.  If all these tests
are infeasible, $d^*$ is at most every attainable profile.
If one test is feasible, no global least profile exists:
any global least profile would also be lexicographically
least, hence would have to equal $d^*$.  The downward
domination by $V_B$ proved above makes these witness tests
global over the original repaired domain as well.  Thus
one lexicographic construction followed by at most one
linear witness test per nonzero coordinate replaces the
finite profile enumeration.

Finally, after existence is certified, let $W^*$ be the
full rational vector space of actual bar-fixed elements
with the prescribed GT support, all lower-coordinate
zero-residue bounds, and ambient degree at most $d^*$.
The current coefficient is not normalized in defining
$W^*$.  The preceding bounded-order membership construction
computes this finite-dimensional space exactly, and
$h_T$ is nonzero on it because a minimum already exists.
Every vector in $W^*\setminus\ker h_T$ normalizes to a
minimum.  Two such vectors have the same normalization
if and only if they span the same rational line.
Since $\mathbb Q$ is infinite, a vector space of dimension
at least two has infinitely many rational lines outside
a proper subspace.  Therefore the final uniqueness test is
simply
\[
                 \text{unique minimizing pair}
                    \quad\Longleftrightarrow\quad
                 \dim_{\mathbb Q}W^*=1.
 \tag{N-77}
\]
This dimension test includes all allowed current exponents
and all possible odd-parity perturbations.

\paragraph{What has now become finite.}
The unbounded scaling and all-degree search in the specified
bar-fixed variant have been reduced to explicit, finite
polynomial-module and rational-linear-algebra operations.
The construction records the chosen label order and initial
GT normalization as inputs and can return a certificate of
failure of existence or uniqueness at a stage.  It does
not resolve N-014 by assigning that chosen order to VIII,
does not modify the refutation of the literal printed
domain in N-012, and does not assert a successful full
rank-four globalization before the finite tests have
actually been run.


\subsection{N-017: an actual minimal block, an order-dependent failure, and a unimodality obstruction}
\label{sec:nonstandard-first-block}

\paragraph{Precise scope of the construction.}
This result executes N-016 on an actual rank-four block.  It retains
the additional bar-fixed condition and rational coefficient field of
N-013, the actual unsaturated order $B_S=\mathbb Q[D]\langle P_i\rangle$
of N-015, and the rational GT matrix units of N-010.  It makes two
different, explicit choices for the bitableau order.  Neither choice
is retroactively attributed to VIII.  The ambient coordinates are
the opposite KL tensor coordinates (N-61).

There is an outer-label repair in which $S'_{22}$ is a minimal
label: use inclusion of the dominance lower sets of maximal
ordinary-weight antichains, keeping distinct equal antichains
incomparable, as in N-014.  The ordinary partition chain at size four is
$4>31>22>211>1111$.  The maximal antichain for $S'_{22}$ is
$\{(22,22)\}$, whereas that for the other minimal label
$\theta S'_{31}$ is $\{(31,211),(211,31)\}$.  These antichains are
incomparable.  The remaining maximal antichains, obtained from the
exact tensor decompositions of N-010, contain one of these two
dominance lower sets.  Thus the first block calculation requires no
choices of scaling in any other outer block.

The two paths of $S'_{22}$ are $N_2-$ and $N_2+$, written simply
$-$ and $+$.  We use $-<+$ on paths.  The first choice is the product
order on the row and column paths:
\[
 (-- )<(-+),(+-)<(++),\qquad (-+)\ \hbox{and}\ (+-)\ 
 \hbox{incomparable}.
 \tag{N-78}
\]
The second choice is its row-first total refinement
$--<-+<+-<++$.  Other outer labels can be placed after this block
in a total refinement when that variant is used.

For the ambient degree comparison, a genuine factor order is
specified by $u\preceq w$ if $u=w$ or $\ell(u)>\ell(w)$, and tensor
coordinates have the product order.  Distinct permutations of equal
length are therefore incomparable.  The calculation also passes,
with the same successful elements and the same failed stage, for
the broader convention comparing numerical length pairs and making
equal length pairs incomparable.  This removes any reliance on that
possible interpretation of VIII's reverse-length wording.  The
comparison of degree tuples is exactly VIII's quantified condition
quoted and analyzed in N-016, rather than a single lexicographic test.

\paragraph{The actual module and its verification.}
Put $L=(D-4)^2(D-2)$ and define
\[
 \begin{aligned}
 e&=D L(D-1)^2,\qquad a=D^2e,\qquad b=D^3 L(D-2)(D-1),\\
 c&=D^3(D-4)^2(D-1)^2,\qquad k=(D-2)/D^2.
 \end{aligned}
 \tag{N-79}
\]
In the GT matrix basis ordered $(+,-)$, the intersection of $B_S$
with this full matrix block has the following actual $S$-basis:
\[
 A=\begin{pmatrix}0&0\\0&a\end{pmatrix},\quad
 B=\begin{pmatrix}0&0\\ b/D&2a/D\end{pmatrix},\quad
 C=\begin{pmatrix}0&c/D\\0&a/(2D)\end{pmatrix},\quad
 E=\begin{pmatrix}e&-2c/D^2\\-b/(2D^2)&-(D+1)e\end{pmatrix}.
 \tag{N-80}
\]
The apparent divisions in these displayed formulas are exact
polynomial divisions.  The rows are, respectively, rows
$113,112,111,110$ of the zero-indexed GT Hermite certificate.

Here are the exact checks establishing that (N-80) describes the
unsaturated intersection.  The N-015 order basis was mapped to the
nine GT blocks and reduced by polynomial Euclidean row operations,
in reverse GT order.  All $114$ pivots are present.  A straight-line
program of $11731$ polynomial row operations gives every new row in
the old order basis; exact polynomial triangular division puts each
old row in the new module.  Thus both module containments hold.
The last four rows have GT support exactly in this block.  Conversely,
any element supported there must have zero coefficients on all earlier
Hermite pivots, successively, so these four rows span the intersection.
The same argument gives the smaller path down-set intersections used
below; in particular $\operatorname{span}_S\{C,A\}$ has no hidden
contribution from $B$, since the excluded $-+$ pivot forces its
coefficient to vanish.

The separate verifier
\nolinkurl{verify_order_gt_hermite.py} does not trust the constructed
word-to-GT evaluation.  It verifies that its proposed map sends the
identity to the block identities and intertwines right multiplication
by each $P_i$ on every one of the $114$ actual order basis rows.
These $342$ identities use the standard-$T$ ambient multiplication
reconstructed in N-015's separate verifier.  Since $1$ generates the
word order under right multiplication, these identities determine
the map on every word.  It then replays the whole Hermite program
jointly in GT and ambient coordinates and checks the reverse
containment by polynomial division.  Consequently neither a
generic projection nor a saturation is used in (N-80).

\paragraph{Exact finite spaces in the globalization problem.}
The first three product-order stages have spaces
\[
 \begin{array}{c|c|c|c}
 T&V&\text{primitive current-line generator}&V/(S t)\\\hline
 --&SA&A&0\\
 -+&S(DB-2A)\oplus\mathbb Q B&DB-2A&\mathbb Q B\\
 +-&S(DC-A/2)\oplus\mathbb Q C&DC-A/2&\mathbb Q C.
 \end{array}
 \tag{N-81}
\]
For example, at $-+$ the previously determined exponent of $--$ is
$-8$, so its lower-coordinate degree bound is $7$.  In $pB+rA$ that
coordinate is $a(2p/D+r)$, with $\deg a=8$.  Writing $p=Du+s$ with
$s\in\mathbb Q$, the bound forces $r=-2u$.  This proves the second
row of (N-81) in the actual order; the third row is identical with
$2$ replaced by $1/2$.

For the last stage put
\[
 T_0=D^2E+2DC+(D/2)B+(D-1)A=aE_{++}.
 \tag{N-82}
\]
The lower exponents for $--,-+,+-$ in the product-order variant
are $-8,-8,-7$.  Directly imposing their exact degree bounds gives
\[
 V=S T_0\oplus
        \operatorname{span}_{\mathbb Q}\{E,C,B,DE+A\}.
 \tag{N-83}
\]
For completeness, write an element as $pE+rC+sB+tA$.  After subtracting
a polynomial multiple of $T_0$, its first coefficient can be taken
as $p=x+wD$.  The $+-$ and $-+$ bounds say respectively
$\deg(Dr-2p)\le1$ and $\deg(Ds-p/2)\le1$, forcing $r,s$ constant.
The $--$ bound says
$\deg(D^2t+Dr/2+2Ds-(D+1)p)\le1$, forcing $t=w$.
This is exactly (N-83).  The coefficient $D^2$ in the primitive
current-line vector is minimal: solving the three lower equations
over $F$ gives coefficient ratios with least common denominator
$D^2$, and their cleared polynomial coefficients have gcd one.

In every row of (N-81), and in (N-83), the cutoff $B$ of (N-73) is
zero.  This is checked by comparing the finitely many ambient
polynomial degrees of the explicitly supplied vectors: for each
nonzero coordinate of the primitive current-line vector, every
complement vector has no greater polynomial degree.  It follows
from N-016 that the spaces to test have dimensions $1,2,2,5$,
respectively.  This is an all-degree reduction, not a degree-limited
search.

The exact linear feasibility tests (N-76)--(N-77) give unique least
profiles in those four finite spaces.  The resulting elements and
exponents are
\[
 \begin{array}{c|c|r}
 T&e_T&a_T\\\hline
 --&X=A&-8\\
 -+&Y=DB-2A&-8\\
 +-&Z=DC-A/2&-7\\
 ++&W=E+C+B/4&-6.
 \end{array}
 \tag{N-84}
\]
The complete inputs, all ambient coordinates, current-coordinate
numerators, quotient bases, and exact outcomes are in
\nolinkurl{bar_fixed_block_rank4_verified.json}; the replay is
\nolinkurl{bar_fixed_block_rank4.py}.  It tests the existence of a
source-comparator failure at every nonzero coordinate of the putative
least profile.  Thus the assertion that these profiles are least
does not follow merely from the lexicographic candidate.

The small linear systems have an additional transparent description.
At $-+$ and $+-$, eliminating the coefficient of
$C_{4321}\otimes C_{4321}$ leaves the displayed pure current-line
vector.  At $++$, the only steps increasing the rank of the greedy
constraints occur at coordinate indices $299,221,272$, representing
\[
 (4321,4321),\qquad(2431,4321),\qquad(3421,3421).
\]
On $(E,C,B,A)$ their normalized ambient coefficient rows are
\[
 L(-1,1/2,2,D),\qquad
 L(-1/2,0,2,D/2),\qquad
 (L/2)(1,0,0,0).
 \tag{N-85}
\]
The first two coordinates are set to zero and the last to degree
at most three.  On (N-83) these conditions leave exactly the line
$\mathbb QW$.  All the remaining source-comparison tests pass.

Full bar-fixed uniqueness also holds.  N-013 excludes an odd current
exponent at a minimum.  Any odd homogeneous perturbation is therefore
$\delta w$ with zero current GT coordinate.  At each of the first
three stages the lower-coordinate bounds force $w=0$.  At the last
stage those bounds permit only $w=\alpha B+\beta C$: the $A$ coefficient
would have degree at least eight whereas its bound is seven.
The first two zero coordinates in (N-85) force
$2\alpha+\beta/2=0$ and $2\alpha=0$, hence $w=0$.
The full minimizing space has rational dimension one, so the
leading-coefficient normalization gives precisely one pair $(a_T,e_T)$.

\paragraph{The total refinement has no fourth global minimum.}
For the row-first total order, the first two elements remain $X,Y$.
At its third stage the lower $-+$ coordinate is also allowed, and
the exact finite space is
\[
 \operatorname{span}_{\mathbb Q}\{C,B,DC-A/2\}.
\]
Its unique full-domain minimum is $C-B/4$ with exponent $-6$.
At the fourth stage the different prior exponent changes the exact
space to
\[
 V=S T_0\oplus
          \operatorname{span}_{\mathbb Q}\{E,B,DE+2C+A\}.
 \tag{N-86}
\]
The cutoff remains zero.  The unique lexicographic minimizing line
has normalized representative
\[
 U=(D+2)E+2C+B/2+A,
 \qquad a_U=-7.
 \tag{N-87}
\]
The admissible element $U'=E+B/2$, with exponent $-6$, contradicts
its being a global least profile: at $(3421,3421)$ its coefficient
is $L/2$, whereas that of $U$ is $(D+2)L/2$.  Its degree is strictly
smaller, and at every strict predecessor its degree is at most that
of $U$.  This last finite assertion is checked directly on the full
ambient coordinates in the certificate; the stricter factor-order
verifier also finds a failure already at $(3412,3412)$.
Thus $d(U)\not\le d(U')$ under the actual source comparator.
Every global least profile would have to be lexicographically least,
so no global minimum exists at this stage.  The distinction between
the two choices of order is therefore mathematically consequential.

\paragraph{Multiplication in the successful block.}
The four elements in (N-84) have matrices
\[
 X=aE_{--},\qquad Y=bE_{-+},\qquad Z=cE_{+-},\qquad
 W=\begin{pmatrix}e&kc\\ kb/4&-e\end{pmatrix}.
 \tag{N-88}
\]
They form an $F$-basis of the matrix block.  Their complete $64$
structure constants, of which $28$ are nonzero, are stored in
\nolinkurl{bar_fixed_block_constants_verified.json} and computed
by \nolinkurl{bar_fixed_block_constants.py}.  Each nonzero constant
has the form
\[
 \varepsilon r\,D^i(D-4)^2(D-2)^j(D-1)^l,
 \quad \varepsilon\in\{1,-1\},\quad r\in\mathbb Q_{>0},
 \quad i\ge1,\quad1\le j\le3,\quad l\ge1.
 \tag{N-89}
\]
Here $D-4=(v-v^{-1})^2$, and $D=q+2+q^{-1}$ with $q=v^2$.
For $j=1,2,3$, respectively, the coefficient lists of
$D(D-2)^j$ from their lowest to highest powers of $q$ are
\[
 (1,2,2,2,1),\quad
 (1,2,3,4,3,2,1),\quad
 (1,2,4,6,6,6,4,2,1).
\]
They are symmetric and unimodal.  The same is true of $D$ and
$D-1$.  Products of nonnegative symmetric unimodal Laurent
polynomials are symmetric unimodal, as follows by expressing each
as a nonnegative sum of centered interval polynomials and multiplying
those intervals.  Equation (N-89) therefore proves the claimed
signed positivity and unimodality of every multiplication constant,
after the factor $(v-v^{-1})^4$ is removed.  Coefficients here are
indexed by every integer $q$ exponent, so parity gaps in $v$ are
not silently discarded from a purported statement about consecutive
$v$ exponents.

\paragraph{The ambient coefficients nevertheless fail unimodality.}
The coefficient of $W$ at the diagonal tensor coordinate
$C_{3421}\otimes C_{3421}$ is exactly
\[
 \frac12(D-4)^2(D-2)
   =(v-v^{-1})^4\,\frac{q+q^{-1}}2.
 \tag{N-90}
\]
The residual coefficient sequence is $(1/2,0,1/2)$, which is not
unimodal even when indexed by consecutive integer powers of $q$.
There is no alternative choice of the vanishing factor making it
nonnegative and unimodal.  Indeed a nonzero Laurent polynomial with
nonnegative coefficients has strictly positive value at $v=1$.
The order of vanishing of (N-90) there is exactly four; hence any
factorization $\pm(v-v^{-1})^m f$ with nonnegative Laurent $f$ forces
$m=4$.  Its sign is then positive, and $f=(q+q^{-1})/2$ is forced.

This is a failure of the structural-coefficient unimodality property
for the explicit repaired variant (N-78), even though the
multiplication constants of its four-element block pass the test
in every product.  It does
not identify this variant with a successfully defined printed $E(4)$,
nor assert that every possible repair has this defect.  The full
$114$-element globalization and its full table have not been
computed by this block calculation.

The obstruction is unchanged if the other KL convention is used
consistently.  The semilinear map $v\mapsto v^{-1}$ and
$T_i\mapsto-T_i$ on both tensor factors fixes every $P_i$ and
$\delta$, and sends $C'_w$ to $(-1)^{\ell(w)}C_w$.
It fixes every element of $B_R^{\mathrm{bar}}$ by
$B_R^{\mathrm{bar}}=B_S\oplus\delta B_S$.
Consequently the two ambient coefficient systems differ only by
the corresponding length signs and bar conjugation; their degrees
and the signed unimodality failure agree.

\paragraph{Source and check boundary.}
The original degree comparator and globalization/scaling domain are
in VIII, section 3.3, printed pages 20--22.  The structural-coefficient
positivity claim is in its subsection ``Positivity in the Kronecker
case.''  The exact word-order and GT inputs are N-010 and N-015,
and the finite all-degree justification is N-016.  The accompanying
computations are exact rational polynomial calculations.  The root
agent independently read this entire proof, both stage helper modules,
the construction and constant-table scripts, and the separate GT-map
verifier, then replayed all three exact computations successfully.
The receipt is \nolinkurl{nonstandard_root_review_N017.md}.


\subsection{N-018: the product-path repair fails in the other minimal block}
\label{sec:nonstandard-second-block-failure}

\paragraph{Statement.}
Use the same explicitly repaired bar-fixed domain, rational GT lines,
outer antichain order and product order on the two paths as in N-017.
The other minimal outer label $\theta S'_{31}$ has five paths, ordered
\[
 N_2-<N_2+<\varepsilon_-<N_1-<N_1+.
\]
Fix the row path $N_2-$ and increase the column path through this list.
The first four stages have unique minima in the full bar-fixed domain.
The fifth has no global minimum.  This failure holds both for the
genuine product of reverse-length factor orders and for the broader
length-stratum ambient order specified in N-017.  Therefore the
product-path repair which succeeds on the whole $S'_{22}$ block does
not give a full rank-four basis.  No statement about every possible
repair of VIII is being made.

\paragraph{The actual five-row module.}
The $25$-dimensional generic block of $\theta S'_{31}$ is isolated in
the actual polynomial order as follows.  The last $29$ rows of the
N-017 full GT Hermite basis span the actual intersection with the sum
of this block and $S'_{22}$.  Put the four $S'_{22}$ coordinates first
and perform polynomial Hermite reduction.  The last $25$ rows then
give the actual intersection with $\theta S'_{31}$ alone.  The
certificate contains $230$ straight-line nodes, exact preimages for
all output rows, and exact polynomial divisions of all input rows.
The pivot argument of N-017 proves both the intersection assertion
and the absence of saturation.  A second such reduction, putting
excluded coordinates first, gives each product-path down-set.

For the five positions in question the down-sets are successive
prefixes in one row, so their actual bases are tails of the isolated
Hermite basis.  Write them, in increasing GT label order, as
$A,B,C,E,F$.  Set
$\rho=D^2(D-4)^2$, $r=D-1$, and $s=D-2$.  Their nonzero GT row entries
are given by the following table, with a common factor $\rho$:
\[
\begin{array}{c|ccccc}
 &F&E&C&B&A\\\hline
 A&0&0&0&0&Ds^2r^2\\
 B&0&0&0&s^2r^2&-s^2r^2/2\\
 C&0&0&Dr^3&0&-3Ds^2r/2\\
 E&0&sr&-sr^2/2&0&D^2s^2/4\\
 F&sr^2&0&r^3&s^2r&-s^2r(2D+1)/2
\end{array}
\tag{N-91}
\]
The column names denote the corresponding GT matrix units, not the
order basis vectors themselves.  All five rows have row path $N_2-$;
their column paths run in the displayed order above.  They are actual
elements of $B_S$, with all $300$ normalized ambient coordinates
stored in the certificates.  In the global reverse GT numbering
their pivots are $109,108,107,106,105$, respectively.  These $A,B,C,E$
are local names for this different block and do not denote the rows
of (N-80).

\paragraph{The first four exact stages.}
The all-degree procedure of N-016 gives
\[
\begin{array}{c|l|r}
\text{current}&\text{unique normalized minimum}&a_T\\\hline
 A&A&-9\\
 B&(D-2)B+A/2&-9\\
 C&C-3B&-8\\
 E&(D^2-3D-6)E+(D-4)C/2-(D-4)A/4&-8
\end{array}
\tag{N-92}
\]
The quotient dimensions $\dim(V/St)$ are $(0,1,2,4)$, the cutoffs of
N-016 are $(0,1,1,1)$, and the finite-space dimensions are $(1,3,4,6)$.
Each minimizing even subspace has rational dimension one,
and every odd minimizing subspace has dimension zero.  Thus these
are full-domain uniqueness statements, not merely minima among the
$\beta$-fixed candidates.

Every entry of this table is determined by exact polynomial
triangular arithmetic and rational linear equations.  The replay
records the primitive current-line coefficient vector, a basis of
$V/(St)$ with actual polynomial lifts, the cutoff, the finite-space
basis, the current GT functional, and the resulting ambient profile.
For each candidate it tests the exact negation of the source
comparison at every nonzero ambient coordinate, as in (N-77).
The complete records are
\nolinkurl{bar_fixed_theta_stages.json}; each actual-module reduction
has its own compressed order certificate.  These finite checks,
together with N-016, verify (N-92) in all degrees.

\paragraph{The fifth-stage quotient can be written explicitly.}
At $F$ the fixed lower exponents of $A,B,C,E$ are $-9,-9,-8,-8$.
Put
\[
 \begin{split}
 T_0={}&D(D-1)F+(1-D)C-DB+(2D-3)A/2,\\
 V={}&S T_0\oplus\operatorname{span}_{\mathbb Q}
       \{F,E,DF+4DE+C,B,2DF+4DE+A\}.
 \end{split}
 \tag{N-93}
\]
This is an exact description of the admissible polynomial space
before current-coordinate normalization.  To verify it directly,
write a general element as $pF+uE+wC+xB+yA$.
The zero-lower-coordinate equations in (N-91) imply $u=0$,
$w=-p/D$ and $x=-p/(D-1)$; hence $D(D-1)$ divides $p$.
The last equation then gives precisely the coefficient of $A$ in
$T_0$.  Its coefficients have gcd one, so it generates the actual
current-line kernel.

Subtract a polynomial multiple of $T_0$ to take
$p=p_0+p_1D$.  The $E$-coordinate degree bound gives $\deg u\le1$.
The $C$ and $B$ bounds force $w,x$ to be constants, and the $A$
bound forces $y$ to be constant.  Write $u=u_0+u_1D$.
The remaining leading-coefficient conditions are
$u_1=2p_1+2w$ and $-4p_1+u_1+4y=0$.
Equivalently $p_1=w+2y$ and $u_1=4w+4y$.
These identities are exactly the five-dimensional complement in
(N-93).  In deriving these equations the lower GT bounds are,
respectively, $8,8,7,7$ for $A,B,C,E$; no denominator or polynomial
divisibility condition has been discarded.

The exact ambient cutoff is $B=1$, so every possible global minimum
lies in the seven-dimensional rational space consisting of the five
displayed complement vectors together with $T_0,DT_0$.
This bound includes arbitrarily high polynomial candidates by the
proof of N-016; it is not a finite experimental degree limit.

\paragraph{A global minimum cannot exist.}
The unique normalized lexicographic candidate in that seven-dimensional
space is
\[
\begin{split}
 U={}&(3D^2-14D-2)F/3-2(7D+9)E/3\\
     &-(3D-7)C/3-(3D-2)B/3+(D-4)A,
          \qquad a_U=-9,\\
 U'={}&(D+1)F+(D+3)E-C/2-B+3A/4,
          \qquad a_{U'}=-8.
\end{split}
\tag{N-94}
\]
Both elements are admissible: their polynomial coefficient vectors
belong to (N-93), their lower bounds hold, and the leading coefficient
of the current GT coordinate is one.  The second is a witness against
the first being a global minimum.  At the ambient coordinate
$z=(3241,4321)$, their coefficients are
\[
 U_z=-\frac{(D-4)^2(3D^2+14D-16)}6,
 \qquad
 U'_z=\frac{(D-4)^2(5D-4)}4.
 \tag{N-95}
\]
The degrees are four and three.  At every strict predecessor $x<z$,
the certificate checks $d(U',x)\le d(U,x)$.
This holds even for the broader numerical length-stratum order,
and consequently for the genuine factor-product order.  Thus there
is no earlier strict improvement allowing $d(U)\le d(U')$ at $z$.
In the factor-product convention there is already a failure at
$(3241,3241)$, where the coefficients are
$-2D(D-4)^2(3D+1)/3$ and $D(D-4)^2$.

The lexicographic minimizing profile is forced: the constraints at
$(4321,4321)$ have rank three in the seven-dimensional space and
set that coefficient to zero; those at $(2431,4321)$ raise the rank
to six and leave coefficient degree two.  The remaining rational
line is the one spanned by $U$.  The exact replay verifies both
these rank statements and all subsequent degree constraints.
Any global least profile must also be lexicographically least in
this linear extension, by N-016.  Since the forced profile fails
the source comparison against $U'$, no global minimum exists.
N-013 and N-016 promote this obstruction from $B_S$ to the full
bar-fixed $B_R$ domain, including all current scaling exponents.

\paragraph{Rational rescaling of GT lines does not repair this failure.}
The conclusion is not an artifact of choosing a particular nonzero
rational representative of each fixed GT line.  More generally,
replace $c'_J$ by $\widetilde c'_J=h_J(v)c'_J$, where
$h_J\in\mathbb Q(v)^\times$.  Let
$s_J=\operatorname{ord}_{q=0}h_J\in\tfrac12\mathbb Z$ and let
$\kappa_J\in\mathbb Q^\times$ be its leading coefficient at zero.
Set
\[
 \widetilde a_J=a_J-s_J.
 \quad\text{Then}\quad
 q^{\widetilde a_J}h_J(v)c'_J
     =u_J(v)q^{a_J}c'_J,
 \qquad u_J\in A_0^\times,\quad u_J(0)=\kappa_J.
 \tag{N-96}
\]
Thus every lower lattice and its bar conjugate are unchanged.
At the current coordinate the target residue is multiplied by
$\kappa_T$, so multiplication of candidates by this nonzero rational
constant is a bijection of the corresponding fibers.  It preserves
bar invariance, actual-order membership, and every ambient degree.
The current exponents are translated bijectively within
$\tfrac12\mathbb Z$.  Inducting over the preceding successful stages
shows that existence, uniqueness and the fifth-stage failure are
unchanged by any such rational normalization of the same GT lines.
This statement concerns the fixed lines and fixed partial orders;
changing either is an additional modification.

\paragraph{Replay and review boundary.}
The construction is
\nolinkurl{bar_fixed_theta_block.py}; the actual isolated order is in
\nolinkurl{bar_fixed_theta_order.json.gz}, followed by five separately
recorded down-set order certificates.  Its inputs are the independently
reviewed N-010/N-015 data and the full GT map/Hermite verification of
N-017.  The root of each new straight-line program is an explicitly
recorded prior actual row.  Both ambient-order conventions are tested
and produce the same first four minima and the same fifth-stage
lexicographic candidate, with the displayed witness valid in both.
The separate certificate-only verifier
\nolinkurl{verify_bar_fixed_theta.py} imports neither the module
constructor nor its incremental equation solver.  It replays the
actual row preimages and triangular containments, then uses exact
Flint matrix ranks to verify every lexicographic bound and every
source-comparator failure test.  Its first four stages have no
failure witnesses; the fifth has $26$ failing coordinates under the
factor-product order and $9$ under the length-stratum order.
The whole-proof and independent-replay statuses are recorded separately
in the branch check ledger.
